\documentclass[11pt,a4paper,oneside]{amsbook}
\usepackage[T1]{fontenc}
\usepackage[utf8]{inputenc}
\usepackage{lmodern}
\usepackage[american]{babel}
\usepackage{microtype}
\usepackage{geometry}
\geometry{margin=1.1in}
\makeatletter
\def\input@path{{../}}
\makeatother
\usepackage{sga3-en}
\usepackage[hidelinks]{hyperref}

\title[SGA~3, Expos\'e XII (English draft)]{Maximal tori, the Weyl group, Cartan subgroups, the reductive center of smooth affine group schemes\\
{\normalsize SGA~3, Expos\'e XII\\English draft}}
\author{A.\ Grothendieck}
\date{}
\hypersetup{
  pdftitle={Maximal tori, the Weyl group, Cartan subgroups, the reductive center of smooth affine group schemes},
  pdfauthor={A. Grothendieck},
  pdfsubject={Unofficial English translation of SGA 3}
}

\begin{document}
\frontmatter
\maketitle

\chapter*{About this draft}
This is an unofficial English translation of Expos\'e~XII, ``Maximal tori, the Weyl group, Cartan subgroups, the reductive center of smooth affine group schemes''
from \textit{Sch\'emas en groupes} (SGA~3).
The source is the recomposition edited by P.~Gille and P.~Polo,
PDF \texttt{Expo12.pdf} (48~pages),
from \url{https://webusers.imj-prg.fr/~patrick.polo/SGA3/}.
The seminar was directed by M.~Demazure and A.~Grothendieck
(IH\'ES, 1962--1964).

Statement numbers follow that PDF. A clear typographical slip is
corrected in the English and marked with a \texttt{\% typo:} comment;
such slips are listed in the README. A slip that would change the
mathematics is left as printed and marked \texttt{\% typo?:}.
Scholarly proofreading remains outstanding.

The translator's contribution is licensed under the MIT License.
The French original remains copyright of its authors and publishers.

\tableofcontents
\mainmatter
\exposetitle{XII}{Maximal tori, the Weyl group, Cartan subgroups, the reductive center of smooth affine group schemes}
\label{XII}
\begin{center}by A.~Grothendieck\end{center}
{\renewcommand{\thefootnote}{0}\footnotetext{xy version of 2/12/08.}}
\setcounter{footnote}{0}

Starting with the present exposé, unlike the preceding exposés, we shall use the known results on the structure of smooth algebraic groups over an algebraically closed field $k$, and above all Borel's theory of affine algebraic groups, presented in the Chevalley Seminar 56/57: ``Classification des groupes de Lie algébriques''\nde{This seminar has now been published, in a form revised by P.~Cartier, as volume~3 of the \textit{Oeuvres de Chevalley} (Springer, 2005).}. Following the usage in the theory of algebraic groups, we refer to this Seminar by the abbreviation BIBLE. For the coming exposés, we shall need in particular the results of BIBLE~4, 5, 6, 7 (the numeral after BIBLE refers to the number of the exposé). It seems moreover that scheme theory brings no notable simplification to Borel's theory as presented in BIBLE. This is why it did not seem useful to reproduce it in the present Seminar, our aim here being to deduce, from known results over algebraically closed fields, analogous results valid over any base prescheme. (This will no longer be so for Chevalley's structure theory of semisimple groups, which, it seems, is treated advantageously ab ovo over an arbitrary base prescheme).

In the oral exposés (which we have followed in Nos.~1 to 4), we had restricted ourselves to group preschemes affine over $S$, relying essentially on the representability theorems of Exp.~XI \No4. In the present notes (\cf Nos.~6 to 8), we show briefly how the affine hypotheses can be eliminated by a simpler method not using the most delicate results of Exposé~XI. For other generalizations of the results contained in the present exposé, see also exposés XV and XVI. (It is obvious that the contents of exposés XI, XII, XV, XVI ought to be completely recast).

\sgasection{1}{Maximal tori}
\numpara{1.0}
First let $G$ be an algebraic group over an \emph{algebraically closed} field $k$. One calls a \emph{maximal torus} of $G$ an algebraic subgroup $T$ of $G$ that is a torus (which means here, since $k$ is algebraically closed, that it is isomorphic to a group of the form $\Gm^r$), and that is maximal for this property. Note that, since $k$ is perfect, $G_{\mathrm{red}}$ is a subgroup of $G$, smooth over $k$, hence is essentially an algebraic group in the sense of BIBLE. Moreover, every reduced subgroup of $G$ (such as a torus) is automatically a subgroup of $G_{\mathrm{red}}$, so the maximal tori of $G$ coincide with the maximal tori of $G_{\mathrm{red}}$. When $G$ is affine, hence $G_{\mathrm{red}}$ affine, a fundamental theorem of Borel tells us that two maximal tori of $G$ are conjugate by an element of $G(k)=G_{\mathrm{red}}(k)$ (BIBLE~6, th.~4 c)), in particular they have the same dimension. We shall call the common dimension of the maximal tori of $G$ the \emph{reductive rank} of $G$. Note moreover that the restriction $G$ affine is unnecessary, as follows from a known theorem of Chevalley asserting that every smooth connected algebraic group over $k$ is an extension of an abelian variety by an affine group; \cf \No6. In Nos.~1 to 4, we shall most often restrict ourselves to group preschemes affine over the base.

Let $G$ be a smooth algebraic group over $k$, and $T$ a maximal torus of $G$; the centralizer $C(T)$ of $T$ in $G$ will be called the \emph{Cartan subgroup} of $G$ associated to $T$. For us this is a group sub-prescheme defined thanks to VIII~6.7, but one will note that since $G$ is smooth over $k$, so is the Cartan subgroup $C$, by XI~5.3, so in this case $C$ is the unique group sub-prescheme of $G$ smooth over $k$ (\ie an algebraic subgroup of $G$ in the sense of BIBLE) such that $C(k)$ is the subgroup of $G(k)$ centralizing $T(k)$, \ie it is essentially the centralizer in the sense of BIBLE. By the conjugacy theorem already cited, the Cartan subgroups associated to the various maximal tori are conjugate to each other, hence have the same dimension. We shall call their common dimension the \emph{nilpotent rank} of $G$; it is equal to that of $G_{\mathrm{red}}$. Let $\rho_r(G)$ and $\rho_n(G)$ be the reductive and nilpotent ranks of $G$; then one has
\[
\rho_r(G)\le\rho_n(G),
\]
and the difference
\[
\rho_u(G)=\rho_n(G)-\rho_r(G)=\dim C/T
\]
could be called, when $G$ is affine, the \emph{unipotent rank} of $G$. When $G$ is smooth, affine and connected, then $C$ is a \emph{nilpotent and connected} algebraic group (BIBLE~6, th.~6 a) and c)), hence (BIBLE~6, th.~2) isomorphic to the product $C_s\times C_u$, where $C_s=T$ is the original maximal torus (which is a fortiori a maximal torus in $C$), and where $C_u$ is a smooth unipotent subgroup, \ie a successive extension of groups isomorphic to $\Ga$ (BIBLE~6, th.~1 cor.~1 and 7, th.~4). One therefore also has, in this case:
\[
\rho_u(G)=\dim C_u.
\]
\begin{remark}{1.1}\label{XII.1.1}
In addition to the three notions of rank that we have just specified for an affine algebraic group, there are two others that are useful, namely the \emph{semisimple rank} $\rho_s(G)$, which is by definition the reductive rank of the quotient $G/R$, where $R$ is the radical of $G$, and finally the \emph{infinitesimal rank} $\rho_i(G)$, which is by definition the nilpotent rank of the Lie algebra of $G$ (which will be defined and studied in the following exposé). We shall not use them in the present exposé. Let us only point out the inequalities:
\[
\rho_s\le\rho_r\le\rho_n\le\rho_i.
\]
\end{remark}
\begin{lemma}{1.2}\label{XII.1.2}
Let $G$ be an algebraic group over the algebraically closed field $k$, $T$ an algebraic subgroup of $G$, $k'$ an algebraically closed extension of $k$, $G'$ and $T'$ the groups deduced from $G,T$ by extension of the base. For $T$ to be a maximal torus of $G$, it is necessary and sufficient that $T'$ be a maximal torus of $G'$.
\end{lemma}
This is an immediate consequence of the ``finite extension principle'' EGA~IV~9.1.1.
\begin{definition}{1.3}\label{XII.1.3}
Let $S$ be a prescheme, $G$ an $S$-group prescheme of finite type, $T$ a group sub-prescheme of $G$. One says that $T$ is a \emph{maximal torus} of $G$ if
\begin{enumerate}[label=\alph*)]
\item $T$ is a torus (IX~1.3) and
\item for every $s\in S$, denoting by $\overline s$ the spectrum of an algebraic closure of $\kappa(s)$, $T_{\overline S}$ is a maximal torus in $G_{\overline s}$.
\end{enumerate}
% typo?: capital S in T_{bar S} kept as printed.
\end{definition}
\begin{remarks}{1.4}\label{XII.1.4}
It follows from 1.2 that when $S$ is the spectrum of an algebraically closed field one recovers the usual definition, and that property 1.3 is stable under every base change. Note that by X~8.8, in definition 1.3 one can replace condition a) by the condition:

 a') $T$ is of finite presentation and flat over $S$.

One should take care that a maximal torus in the sense of definition 1.3 is indeed maximal in the set of subtori of $G$ (as follows at once from IX~2.9), but that the converse can be exact, the base $S$ being, say, connected, only if $G$ actually admits a maximal torus in the sense of 1.3, which is not the case in general (even if $S=\Spec(\ZZ)$ and if $G$ is ``semisimple''; see also 1.6). We shall see, however, in XIV that the converse is true when $S$ is artinian, or when $S$ is a local scheme and $G$ is ``reductive'': in this case, every torus of $G$ is contained in a maximal torus.
\end{remarks}
\begin{definition}{1.5}\label{XII.1.5}
Let $G$ be an algebraic group over a field $k$. One calls the \emph{reductive rank} (resp. \emph{nilpotent rank}, resp. \emph{unipotent rank}, etc.) of $G$ the reductive rank (resp.\ $\ldots$) of $G_{\overline k}$, where $\overline k$ is an algebraic closure of $k$.
\end{definition}
One will note that, by 1.2 and the commutation of the formation of $\uCentr_G(T)$ with extension of the base, the notions of rank introduced in 1.5 are \emph{invariant under extension of the base field}; on the other hand, for algebraically closed $k$, they coincide with those introduced at the beginning of the present number.
\begin{remark}{1.6}\label{XII.1.6}
It is not difficult to construct an affine smooth group scheme $G$ over the spectrum $S$ of a discrete valuation ring, whose generic fiber is isomorphic to $\Gm$, and whose special fiber is isomorphic to $\Ga$. Such a $G$ contains no torus except the trivial torus $T$ (reduced to the identity subgroup), which is obviously not a maximal torus. More precisely, in the special fiber $G_0=\mathbf G_{\mathrm a,k}$, $T_0$ is indeed a maximal torus, but in the generic fiber $G_1=\mathbf G_{\mathrm m,K}$, $T_1$ is no longer a maximal torus ($k$ is the residue field, $K$ the fraction field of the valuation ring considered). One also sees from this example that the reductive rank of $G_s$ ($s\in S$) is not a continuous function of $s$. One nevertheless has the following results:
\end{remark}
\begin{theorem}{1.7}\label{XII.1.7}
Let $S$ be a prescheme, $G$ an $S$-group prescheme affine and smooth over $S$. For every $s\in S$, consider $\rho_r(s)=\rho_r(G_s)$ and $\rho_n(s)=\rho_n(G_s)$, the reductive and nilpotent ranks of $G_s$ (1.5). With this notation, one has the following:
\begin{enumerate}[label=\alph*)]
\item The function $\rho_r$ on $S$ is lower semicontinuous, the function $\rho_n$ on $S$ is upper semicontinuous, hence the function $\rho_u=\rho_n-\rho_r$ is upper semicontinuous.
\item The following conditions (stable under arbitrary base change) are equivalent:
\begin{enumeratei}
\item The function $\rho_r$ on $S$ is locally constant.
\item There exists locally, in the sense of the étale topology, a maximal torus in $G$.
\item[(ii bis)] There exists locally, in the sense of the faithfully flat quasi-compact topology, a maximal torus in $G$.
\end{enumeratei}
\item Let $T_1,T_2$ be two maximal tori in $G$ (which implies that one is under the conditions considered in b)). Then $T_1,T_2$ are locally conjugate in the sense of the étale topology, \ie there exists an étale surjective morphism $S'\to S$ such that the subgroups $(T_1)_{S'}$ and $(T_2)_{S'}$ of $G_{S'}$ are conjugate by a section of $G_{S'}$ over $S'$.
\item Under the conditions of b), \ie when $\rho_r$ is locally constant, so is $\rho_n$ (hence also $\rho_u=\rho_n-\rho_r$).
\end{enumerate}
\end{theorem}
\begin{proof}
a) Note that for every morphism $S'\to S$, if $G'=G\times_S S'$, the functions $\rho'_r$, etc., on $S'$, defined in terms of $G'$ as $\rho_r$, etc., in terms of $G$, are obtained simply by composing the latter with $S'\to S$. When $S'\to S$ is faithfully flat quasi-compact, it follows that $\rho'$ is continuous, resp. upper semicontinuous, resp. lower semicontinuous, if and only if $\rho$ is, the Zariski topology of $S$ being indeed the quotient of that of $S'$ (SGA~1 VIII~4.3). Consequently, the assertions of a) are local for the faithfully flat quasi-compact topology. Let then $s\in S$; one wants to show that the set $U$ of $t\in S$ such that $\rho_r(t)\ge\rho_r(s)$ (resp. $\rho_n(t)\le\rho_n(s)$) is a neighborhood of $s$. By the finite extension principle, there exists a finite extension $k$ of $\kappa(s)$ such that $G_k$ admits a maximal torus. There then exists an open neighborhood $U$ of $s$, and a finite surjective locally free morphism $S'\to U$, such that the fiber $S'_s$ is $\kappa(s)$-isomorphic to $\Spec(k)$ (\cf EGA~III~10.3.2, where the noetherian hypothesis is manifestly unnecessary). Since $S'\to S$ is an open morphism (SGA~1 IV~6.6), one is reduced to the case where $S'=S$, \ie to the case where there exists a maximal torus $T_s$ in $G_s$. In addition, thanks to XI~5.8 a), on replacing $S$ further by an $S'$ étale over $S$ and endowed with a point $s'$ above $s$, one can suppose that $T_s$ is the fiber at $s$ of a subtorus $T$ of $G$. Then for every $t\in S$, $\rho_r(t)=\rho_r(G_t)\ge\dim T_t=\dim T_s=\rho_r(G_s)=\rho_r(s)$, which proves that $\rho_r$ is lower semicontinuous. On the other hand, by XI~5.3, the functor
\[
C=\uCentr_G(T)
\]
is representable by $C=\Centr_G(T)$\nde{Modification made to introduce the prescheme $C=\Centr_G(T)$ representing the functor $\uCentr_G(T)$.}, a closed group sub-prescheme of $G$ smooth over $S$. Thus there exists a neighborhood $U$ of $s$ such that $t\in U$ implies $\dim C_t=\dim C_s=\rho_n(s)$. The upper semicontinuity of $\rho_n$ is then a consequence of the relation
\[
\rho_n(t)\le\dim C_t\quad\text{for every }t\in S,
\]
which is itself contained in the following purely geometric lemma:
\begin{lemma}{1.8}\label{XII.1.8}
Let $k$ be a field, $G$ a smooth affine algebraic group over $k$, $T$ a torus in $G$, $C$ its centralizer; then one has
\[
\rho_n(G)\le\dim C.
\]
\end{lemma}
Indeed, one can suppose $k$ algebraically closed, so that $T$ is contained in a maximal torus $T'$. Let $C'$ be the centralizer of the latter; then one has $C'\subset C$, hence $\dim C'\le\dim C$. \hfill Q.E.D.

b) If $\rho_r$ is locally constant, then for every torus $T$ in $G$, and every $s\in S$, if $T_s$ is a maximal torus in $G_s$, then there exists an open neighborhood $U$ of $s$ such that $T|_U$ is a maximal torus in $G|_U$. Using now the argument of a), one sees that (i) $\Rightarrow$ (ii bis). On the other hand (ii bis) $\Rightarrow$ (i), for when $G$ admits a maximal torus $T$, then it is obvious that $\rho_r(s)=\dim T_s$ is a locally constant function of $s$, but we pointed out at the beginning of the proof of a) that the question of the continuity of $\rho_r$ was local for the faithfully flat topology. Thus (i) $\Leftrightarrow$ (ii bis), and obviously (ii) $\Rightarrow$ (ii bis); it remains to prove the converse implication (i) $\Rightarrow$ (ii). For this, introduce the functor $F$ of XI~4.1\nde{This is the functor of subgroups of multiplicative type of $G$.}, which is thus a smooth separated prescheme over $S$, and consider the subfunctor $\mathscr T$ of $F$, whose value at an $S'$ over $S$ is the set of maximal tori in $G_{S'}$. Using the preceding remark that, assuming (i), a torus in $G_{S'}$ that is maximal in the fiber of a point $s\in S'$ is maximal over an open neighborhood of $s$, one sees that $\mathscr T$ is representable by an open sub-prescheme of $F$, and consequently is smooth and separated over $S$. Since the structural morphism $\mathscr T\to S$ is obviously surjective, it therefore admits a section locally in the sense of the étale topology by XI~1.10, which proves (i) $\Rightarrow$ (ii).

c) This is an immediate consequence of XI~5.4 bis, taking account of Borel's conjugacy theorem recalled at the beginning of the number.

d) Assuming the remark at the beginning of the proof in a), and taking account of b), one can suppose that there exists a maximal torus $T$ in $G$. If $C$ is its centralizer, then $C$ is representable and is smooth over $S$ by XI~5.3, so the function $s\mto\rho_n(s)=\dim C_s$ is indeed locally constant.

The proof of 1.7 is complete. We shall refer to the conditions considered in 1.7 b) by saying that, in this case, $G$ has \emph{locally constant reductive rank}. Note:
\end{proof}
\begin{corollary}{1.9}\label{XII.1.9}
Let $G$ be as in 1.7, and let $s\in S$ be such that $\rho_u(s)=0$, \ie $\rho_r(s)=\rho_n(s)$ (\ie the Cartan subgroups of $G_{\overline k}$ are tori, where $k$ is an algebraic closure of $\kappa(s)$). Then there exists an open neighborhood $U$ of $s$ on which $\rho_r$ and $\rho_n$ are constant; in particular, for every $t\in U$, the unipotent rank $\rho_u(t)$ of $G_t$ is zero.
% typo?: G_{bar k}, although k already denotes an algebraic closure, kept as printed.
\end{corollary}
Indeed, this follows immediately from 1.7 a) and the inequality $\rho_r(t)\le\rho_n(t)$ for every $t\in S$.

Note also that we proved, at the same time as b), the
\begin{corollary}{1.10}\label{XII.1.10}
Let $G$ be as in 1.7, and suppose $G$ of locally constant reductive rank. Consider the functor
\[
\mathscr T:\Sch^\circ\to\Ens,
\]
such that for every $S'$ over $S$, one has
\[
\mathscr T(S')=\text{set of maximal tori of }G_{S'}.
\]
Then $\mathscr T$ is representable by a prescheme smooth, separated and of finite type over $S$.
\end{corollary}
It remains to verify that $\mathscr T$ is of finite type over $S$. Since the question is local for the faithfully flat quasi-compact topology, one can suppose that $G$ admits a maximal torus $T$. By XI~5.3 bis and the bis version of 5.5, $\uNorm_G(T)$ and $G/\uNorm_G(T)$ are representable by preschemes $\Norm_G(T)$ and $G/\Norm_G(T)$, and $\mathscr T$ is isomorphic to $G/\Norm_G(T)$\nde{Modification made to introduce the preschemes $\Norm_G(T)$ and $G/\Norm_G(T)$.}. Since the morphism $G\to\mathscr T$ defined by $g\mto\operatorname{int}(g)(T)$ is surjective, and $G$ quasi-compact over $S$, so is $\mathscr T$, which completes the proof.
\begin{remarks}{1.11}\label{XII.1.11}
a) The prescheme $\mathscr T$ of 1.10 will be called, as is natural, the \emph{prescheme of maximal tori} of $G$. One will see in \No5 that it is in fact \emph{affine} over $S$. One can verify this directly when $G$ is isomorphic to a closed group sub-prescheme of a prescheme of the form $\GL(n)$, using XI~4.6 and observing that, by the argument of the proof of 1.7 b), $\mathscr T$ is identified with a sub-prescheme, both \emph{open and closed}, of the prescheme $F$ representing the subgroups of multiplicative type of $G$.

b) It is possible to give a proof of 1.7, hence of 1.9, not using the delicate results of XI, but only the easy results XI~3.12 and 6.2, working solely with groups of multiplicative type finite over the base (morally, the groups ${}_nT$ where $T$ is a maximal torus), compare \No7. The same remark holds for the proof of 1.10.
\end{remarks}

We finish this number by giving examples where there exists a unique maximal torus.
\begin{proposition}{1.12}\label{XII.1.12}
Let $G$ be an $S$-group prescheme of multiplicative type and of finite type. Then $G$ admits a unique maximal torus, and every torus in $G$ is contained in this maximal torus.
\end{proposition}

Uniqueness follows obviously from this last assertion, which characterizes the maximal torus as the largest subtorus of $G$. From uniqueness it follows that the existence question is local for the faithfully flat quasi-compact topology, which allows us to suppose $G$ diagonalizable, \ie of the form $D_S(M)$, $M$ a commutative group of finite type. Let $M_0$ be the quotient of $M$ by its torsion subgroup; I say that the torus $T=D_S(M_0)$ in $G$ is a maximal torus and a largest subtorus. Indeed, a subtorus $T'$ of $G$ is locally diagonalizable for the fpqc topology, so to prove $T'\subset T$, one can suppose $T'$ diagonalizable, hence of the form $D_S(N)$, where $N$ is a free quotient of $M$ (VIII~1.4 and 3.2 b)), hence $N$ is a quotient of $M_0$, so $T'\subset T$. Since the construction of $T$ as $D_S(M_0)$ is compatible with every extension of the base, this shows at the same time that $T$ is a maximal torus of $G$, and completes the proof. In the case where $G$ is smooth over $S$, one can generalize 1.12:
\begin{proposition}{1.13}\label{XII.1.13}
Let $G$ be an $S$-group prescheme of finite presentation over $S$. Suppose that $G$ admits locally for the faithfully flat quasi-compact topology a central maximal torus. Then $G$ admits (globally) a unique maximal torus\nde{Which is central!}, and it is the largest subtorus of $G$.
\end{proposition}
Here again, uniqueness is trivial from the last assertion, and makes existence a local question for fpqc, which allows us to suppose that $G$ admits a central maximal torus $T$. Let us show that every torus $R$ of $G$ is contained in $T$.

This follows from the
\begin{lemma}{1.14}\label{XII.1.14}
Let $G$ be an $S$-group prescheme of finite presentation over $S$, $T$ a maximal torus of $G$, $R$ a subtorus of $G$, and suppose that $T$ and $R$ commute. Then $R\subset T$.
\end{lemma}
Indeed, since $R$ and $T$ commute, the morphism $R\times_S T\to G$ defined by $(r,t)\mto r\cdot t$ is a group homomorphism, so by IX~6.8 it admits an image subgroup $T'$ in $G$, which is a group of multiplicative type quotient of $R\times_S T$, hence a torus, obviously containing $T$. Since $T$ is a maximal torus, one has $T'=T$ (1.4), so $R\subset T$, which proves 1.14, hence 1.13. In particular, using 1.7 b):
\begin{corollary}{1.15}\label{XII.1.15}
Let $G$ be a commutative $S$-group prescheme, smooth and affine over $S$, with locally constant reductive rank. Then $G$ admits a unique maximal torus, and the latter contains every subtorus of $G$.
\end{corollary}
\begin{corollary}{1.16}\label{XII.1.16}
Let $G$ be an $S$-group prescheme smooth and affine over $S$. Suppose that for every $s\in S$, denoting by $\overline s$ the spectrum of an algebraic closure $k$ of $\kappa(s)$, the geometric fiber $G_{\overline s}$ is a connected and nilpotent algebraic group (in the sense of BIBLE, \ie the group $G(k)$ of its points with values in $k$ is nilpotent). Suppose in addition the reductive rank of $G$ locally constant. Then $G$ admits a unique maximal torus $T$; in addition $T$ is central and is the largest subtorus of $G$.
\end{corollary}
By 1.7 b), $G$ admits locally for fpqc a maximal torus, and by 1.13 one is reduced to proving that a maximal torus of $G$ is central. By IX~5.6 b), one is reduced to the case where $S$ is the spectrum of a field, which one can suppose algebraically closed. Then $T(k)$ is in the center of $G(k)$ by BIBLE~6 th.~2, which implies that $T$ is in the center of $G$, for $\Centr_G(T)$ is a closed subscheme of $G$ (VIII~6.6), containing the points of $G(k)$, hence is identical to $G$ by the Nullstellensatz ($G$ being reduced).

Finally, for later reference, let us point out the following trivial proposition (which we have already used implicitly):
\begin{proposition}{1.17}\label{XII.1.17}
Let $G\supset H\supset T$ be group preschemes of finite type over $S$. Equivalent conditions:
\begin{enumeratei}
\item $T$ is a maximal torus of $G$.
\item $T$ is a maximal torus of $H$, and for every $s\in S$, one has equality of the reductive ranks $\rho_r(G_s)=\rho_r(H_s)$.
\end{enumeratei}
\end{proposition}

\sgasection{2}{The Weyl group}
First let $k$ be an algebraically closed field, and let $G$ be an algebraic group over $k$, smooth and affine over $k$. If $T$ is a maximal torus, $C$ its centralizer and $N$ its normalizer, then by XI~5.9 these are smooth closed subgroups of $G$, and $C$ is an open subgroup of $N$, so $W=N/C$ is a finite étale group over $k$, hence determined by the group $W(k)$ of its points with values in $k$, as $W=W(k)_k$ (the constant group defined by the ordinary finite group $W(k)$). The finite group $W(k)$ will be called the \emph{geometric Weyl group} (or simply Weyl group if no confusion is to be feared) \emph{of $G$ relative to $T$}. By Borel's conjugacy theorem, the Weyl groups relative to the various maximal tori are isomorphic to each other, which is why one sometimes speaks of ``the'' Weyl group of $G$, without specifying a maximal torus. Since the formation of $C,N$, and $N/C$ commutes with every extension of the base, one sees that if $k'$ is an algebraically closed extension of $k$, the geometric Weyl group of $G_{k'}$ relative to $T_{k'}$ is canonically isomorphic to that of $G$ relative to $T$; consequently ``the'' geometric Weyl group of $G$ (which is properly speaking an isomorphism class of ordinary finite groups) coincides with that of $G'$.

This allows one, when $G$ is a smooth affine algebraic group over an arbitrary field $k$, to speak of the \emph{geometric Weyl group} of $G$ as the isomorphism class of the geometric Weyl group of $G_{k'}$, where $k'$ is any algebraically closed extension of $k$. When $G$ admits the maximal torus $T$, one can obviously form as above $C=\Centr_G(T)$, $N=\Norm_G(T)$, $W=N/C$, which is a finite étale group over $k$, called the \emph{Weyl group of $G$ relative to $T$}; the geometric Weyl group is then nothing other than the class of the group of points of $W$ with values in any algebraically closed extension $k'$ of $k$. Here, knowledge of the geometric Weyl group $W(k')$ obviously no longer suffices in general to reconstruct the algebraic group $W$: one must know in addition the action on $W(k')$ of the Galois group of the separable algebraic closure of $k$ in $k'$.

When finally $G$ is a group prescheme over an arbitrary base $S$, $G$ smooth and affine over $S$, and if $T$ is a maximal torus of $G$, then XI~5.9 again allows us to form the group
\[
W(T)=N(T)/C(T),
\]
which is an $S$-group prescheme étale, separated and quasi-finite over $S$. Its geometric fibers (relative to the algebraic closures of the residue fields $\kappa(s)$, $s\in S$) are the geometric Weyl groups of the fibers $G_s$. Consequently XI~5.10 gives us information on the variation of these groups with $s\in S$. We can make this information more precise and generalize it in the following way:
\begin{theorem}{2.1}\label{XII.2.1}
Let $S$ be a prescheme, $G$ an $S$-group prescheme smooth and affine over $S$. For every $s\in S$, let $w(s)$ be the geometric Weyl group of $G_s$, which is thus an isomorphism class of finite groups. In the set $E$ of such classes, introduce the following preorder relation: $w\le w'$ if and only if $w$ and $w'$ are represented by finite groups $W$ and $W'$ such that $W$ is isomorphic to a quotient of a subgroup of $W'$. With this convention:
\begin{enumerate}[label=\alph*)]
\item The function $s\mto w(s)$ from $S$ to $E$ is lower semicontinuous.
\item Suppose that the reductive rank of $G$ is locally constant. Then the following conditions are equivalent:
\begin{enumeratei}
\item The function $s\mto w(s)$ is locally constant.
\item[(i bis)] The function $s\mto\operatorname{card}w(s)$ is locally constant.
\item There exists, locally for the étale topology (or merely for the fpqc topology), a maximal torus $T$ such that $W(T)$ is finite over its base.
\item[(ii bis)] For every $S'$ over $S$, and every maximal torus $T$ of $G_{S'}$, the associated Weyl group $W(T)$ is finite over $S'$.
\end{enumeratei}
\end{enumerate}
\end{theorem}
\begin{proof}
a) Proceeding as in 1.7 a), to prove that for every $s\in S$ there exists an open neighborhood $U$ of $s$ such that $t\in U$ implies $w(t)\ge w(s)$, one is reduced to the case where there exists a torus $R$ in $G$ such that $R_s$ is a maximal torus in $G_s$. Let
\[
W(R)=\uNorm_G(R)/\uCentr_G(R)
\]
as in XI~5.9; it is a group prescheme étale, separated, quasi-finite over $S$. For every $t\in S$, let $w'(t)\in E$ be its geometric fiber at $t$. Since $R_s$ is a maximal torus in $G_s$, and the formation of $\Norm$, $\Centr$, $\Norm/\Centr$ is compatible with every extension of the base, in particular passage to fibers, one sees that one has
\[
w(s)=w'(s);
\]
I say in addition that, for $t$ near $s$, one will have the inequalities
\[
w(t)\ge w'(t)\quad\text{and}\quad w'(t)\ge w'(s),
\]
which will suffice to establish a). These two inequalities are contained in the following two lemmas:
\begin{lemma}{2.2}\label{XII.2.2}
Let $S$ be a prescheme, $W$ an $S$-group prescheme that is étale, separated and quasi-finite over $S$. For every $s\in S$, let $f(s)$ be the class of the geometric fiber of $W$ at $s$, which is an element of the ordered set $E$ of isomorphism classes of finite groups introduced in 2.1. Then the function $f:S\to E$ is lower semicontinuous. For it to be constant near $s\in S$, it is necessary and sufficient that the same hold for the function $s\mto\operatorname{card}f(s)$, and for this it is necessary and sufficient that $W$ be finite over $S$ above an open neighborhood $U$ of $s$.
\end{lemma}
This result is a refinement, in terms of groups, of that invoked in the proof of XI~5.10; we restrict ourselves to a sketch of the proof (which is of the most standard type). One is reduced as usual to the case of affine noetherian $S$. One sees immediately that the function $f$ is constructible (EGA~$0_{\mathrm{III}}$~9.3.1 and 9.3.2, and the sorites of EGA~IV~9), and by EGA~$0_{\mathrm{III}}$~9.3.4, for semicontinuity one is reduced to proving that if $t$ is a generalization of $s$, one has $f(t)\ge f(s)$. This reduces us, thanks to EGA~II~7.1.7, to the case where $S$ is the spectrum of a discrete valuation ring, which one can suppose complete with algebraically closed residue field. But then, $s$ denoting the closed point of $S$, since $G$ is étale and separated over $S$, it contains a subscheme $G'$ both open and closed, finite over $S$, such that $G'_s=G_s$ (EGA~II~6.2.6), and one sees immediately that $G'$ is here a subgroup of $G$. Moreover, since $G'$ is finite étale over $S=\Spec(V)$, $V$ complete with algebraically closed residue field, it is a constant group, hence of the form $A_S$, where $A=G'(\kappa(s))=G(\kappa(s))$ has class $f(s)$. If $B$ is the geometric fiber of $G$ at the generic point $t$ of $S$, one therefore has a canonical monomorphism $A\to B$, which proves $f(s)\le f(t)$. (N.B. This proof in fact proves semicontinuity for an order relation on $E$ finer than that indicated in 2.1). The fact that $f$ is continuous at $s$ if and only if $s\mto\operatorname{card}f(s)$ is follows from the fact that for $w,w'\in E$, the relations $w\le w'$ and $\operatorname{card}w=\operatorname{card}w'$ imply $w=w'$. The fact that this condition is equivalent to the finiteness of $W$ over a neighborhood of $s$ is then independent of the group structure on $W$, and was pointed out after XI~5.10; moreover its proof is easily given by the preceding arguments, using the valuative criterion of properness EGA~II~7.3.8.
% typo?: G used instead of W throughout this proof; retained.
\begin{lemma}{2.3}\label{XII.2.3}
Let $G$ be a smooth affine algebraic group over an algebraically closed field $k$, $R\subset T$ two subtori, $W(R)$ and $W(T)$ the two finite groups associated as in XI~5.9, quotients of the normalizer by the centralizer. Then $W(R)$ is isomorphic to a quotient of a subgroup of $W(T)$.
% typo?: T not specified to be maximal in the statement, though used in proof; kept.
\end{lemma}
Indeed, consider the diagram
\[
\xymatrix{T\ar@{^{(}->}[r]&C(T)\ar@{^{(}->}[r]\ar@{^{(}->}[d]&C(R)\ar@{^{(}->}[d]\\
&N(T)\cap N(R)\ar@{^{(}->}[r]\ar@{^{(}->}[d]&N(R)\\
&N(T)&}
\]
then $(N(T)\cap N(R))/C(T)$ is a subgroup of $W(T)=N(T)/C(T)$, and one has an obvious homomorphism:
\[
(N(T)\cap N(R))/C(T)\to W(R)=N(R)/C(R),
\]
and everything amounts to proving that the latter is surjective, hence that for every point $g$ of $N(R)$ with values in $k$, there exists a point $c$ of $C(R)$ with values in $k$ such that $cg$ normalizes $T$, \ie such that
\[
\operatorname{int}(c)(\operatorname{int}(g)T)=T.
\]
Now for this, it suffices to note that $\operatorname{int}(g)T$ is a torus of $N(R)$, hence of $C(R)$ (which is an open subgroup of it). Then $T$ and $\operatorname{int}(g)T$ are maximal tori of $C(R)$, since they are maximal in $G$, and one concludes by Borel's conjugacy theorem.

This proves 2.3 and thereby 2.1 a).

b) We have already pointed out that (i) and (i bis) are trivially equivalent; they imply (ii bis) by the converse of 2.2 or XI~5.10, as one chooses; (ii bis) $\Rightarrow$ (ii) thanks to 1.7 b), and finally (ii) $\Rightarrow$ (i bis), for one sees as in 1.7 a) that condition (i) is local for fpqc, which allows us to suppose that $G$ admits a maximal torus $T$ such that $W(T)$ is finite over $S$, and one again concludes by XI~5.10. This completes the proof of 2.1.
\end{proof}

\sgasection{3}{Cartan subgroups}
\begin{definition}{3.1}\label{XII.3.1}
Let $G$ be a smooth group prescheme of finite type over the prescheme $S$. One calls a \emph{Cartan subgroup} of $G$ a group sub-prescheme $C$ of $G$, smooth over $S$, such that for every $s\in S$, denoting by $\overline s$ the spectrum of an algebraic closure of $\kappa(s)$, $C_{\overline s}$ is a Cartan subgroup (\cf \No1) of $G_{\overline s}$.
\end{definition}
It is immediate that if $C$ is a Cartan subgroup of $G$, then for every $S'$ over $S$, $C_{S'}$ is a Cartan subgroup of $G_{S'}$. If $S$ is the spectrum of an algebraically closed field, one recovers the notion recalled in \No1. Finally, one verifies at once that the fact that a group sub-prescheme $C$ of $G$ is a Cartan subgroup is \emph{of a local nature} for the faithfully flat quasi-compact topology.
\begin{theorem}{3.2}\label{XII.3.2}
Let $G$ over $S$ be as in 3.1, and suppose that $G$ is affine over $S$ and of locally constant reductive rank. Then the map
\[
T\mto\Centr_G(T)
\]
induces a bijection from the set of maximal tori of $G$ to the set of Cartan subgroups{\renewcommand{\thefootnote}{*}\footnote{The proof given here in fact proves 3.2 only for closed Cartan subgroups of $G$. However, 7.1 a) establishes 3.2 in the form stated, and implies that the Cartan subgroups of $G$ are closed.}\addtocounter{footnote}{-1}} of $G$. If $C$ corresponds to $T$, then $T$ is the unique maximal torus of $C$.
\end{theorem}
If $T$ is a maximal torus of $G$, the functor $\uCentr_G(T)$ is representable, by XI~5.3 or 6.2, by a closed smooth sub-prescheme of $G$, $C=\Centr_G(T)$, and it follows from the definitions that $C$ is a Cartan subgroup of $G$. In addition, $T$ is obviously a maximal torus of $C$, and being central, is the unique maximal torus of $C$ (1.13). Thus the map $T\mto\Centr_G(T)$ from the set of maximal tori to the set of Cartan subgroups is injective; it remains to prove that it is surjective. Thus let $C$ be a Cartan subgroup of $G$; let us prove that it is of the form $\Centr_G(T)$, for $T$ a maximal torus of $G$. For this it suffices to find a maximal torus $T$ of $C$, for then $T$ will be a maximal torus of $G$ (for, for every $s\in S$, $G_s$ and $C_s$ have the same reductive rank). By IX~5.6 b), $T$ is in the center of $C$, so $C\subset C'=\Centr_G(T)$, and then $C$ is a smooth subgroup of the smooth group $C'$ over $S'$, coinciding with $C'$ fiber by fiber, whence $C=C'$. Now, since $G$ is of locally constant reductive rank, so is $C$, hence $C$ admits a maximal torus \emph{locally} for the étale topology by 1.7 b), and since this torus is central by the preceding argument, it follows by 1.13 that $C$ does indeed admit a maximal torus, which completes the proof.
% typo?: base S' in the middle of the proof kept as printed.
\begin{corollary}{3.3}\label{XII.3.3}
Let $G$ be a group prescheme smooth and affine over $S$ of locally constant reductive rank, let $\mathscr C:\Sch^\circ_{/S}\to\Ens$ be the functor defined by
\[
\mathscr C(S')=\text{set of Cartan subgroups of }G_{S'}.
\]
Then the functor $\mathscr C$ is isomorphic to the functor $\mathscr T$ of 1.10, hence is representable by the same prescheme smooth, separated and of finite type over $S$.
\end{corollary}
\begin{corollary}{3.4}\label{XII.3.4}
Under the conditions of 3.2, if $C=\Centr_G(T)$, one has
\[
\uNorm_G(C)=\uNorm_G(T).
\]
\end{corollary}
\begin{remark}{3.5}\label{XII.3.5}
When $G$ is not of locally constant reductive rank, it is nevertheless possible that $G$ admits Cartan subgroups (for example if $G$ has connected nilpotent fibers, $G$ is a Cartan subgroup of itself, but is not necessarily of locally constant reductive rank, \cf 1.6). In XV, we develop the theory of Cartan subgroups without supposing $G$ affine over $S$ or of locally constant reductive rank, using the theory of regular elements of the next exposé. See also Nos.~6 and 7 for the elimination of the affine hypothesis.
\end{remark}

\sgasection{4}{The reductive center}
\begin{definition}{4.1}\label{XII.4.1}
Let $S$ be a prescheme, $G$ an $S$-group of finite presentation over $S$ with affine fibers, $Z$ a group sub-prescheme of $G$. One says that $Z$ is a \emph{reductive center} of $G$ if
\begin{enumeratei}
\item $Z$ is central, and of multiplicative type, and
\item for every base change $S'\to S$, and every central homomorphism $u:H\to G_{S'}$, where $H$ is a group of multiplicative type and of finite type over $S$, $u$ factors through $Z_{S'}$.
\end{enumeratei}
(For a variant of this notion when the fibers of $G$ are not supposed affine, \cf 8.6).
% typo?: H over S rather than S' kept as printed.
\end{definition}

Note that such a $Z$ is necessarily of finite type over $S$ (since its fibers are), hence is uniquely determined as the largest central subgroup of multiplicative type of $G$. It is easy to give examples where $G$ (smooth and affine over $S$) admits a largest central subgroup $Z$ of multiplicative type, but where $Z$ is not a reductive center (\ie $G$ admits no reductive center), \cf for example 1.6; it nevertheless follows easily from IX~6.8 that a subgroup $Z$ of $G$ is a reductive center of $G$ if and only if it is a largest central subgroup of multiplicative type, and retains this property under every base change; see also 4.3 below.

It is obvious from 4.1 that if $Z$ is the reductive center of $G$, then for every base change $S'\to S$, $Z_{S'}$ is the reductive center of $G_{S'}$. From this, and from the uniqueness of the reductive center, there follows, thanks to faithfully flat quasi-compact descent theory (SGA~1 VIII):
\begin{proposition}{4.2}\label{XII.4.2}
Let $G$ be an $S$-group prescheme of finite presentation over $S$ with affine fibers. If $G$ admits a reductive center locally for the faithfully flat quasi-compact topology, it admits a reductive center. For $Z$ to be a reductive center of $G$, it is necessary and sufficient that it be so locally for the fpqc topology.
\end{proposition}
Note also:
\begin{proposition}{4.3}\label{XII.4.3}
Let $G$ be an $S$-group prescheme of finite presentation with affine fibers, $Z$ a group sub-prescheme. For $Z$ to be a reductive center of $G$, it is necessary and sufficient that it be of multiplicative type, and that for every $s\in S$, $Z_s$ be a reductive center of $G_s$.
\end{proposition}
Indeed, it follows first from IX~5.6 b) that $Z$ is then central. Since the conditions considered are stable under base change, it remains to prove that every central homomorphism $u:H\to G$, with $H$ of multiplicative type and of finite type, factors through $Z$. Now, since $Z$ is central, $u$ and the canonical immersion $v:Z\to G$ define a group homomorphism
\[
w:H\times_S Z\to G.
\]
By IX~6.8 the latter admits an image group $K$, which is a subgroup of multiplicative type of $G$, and everything amounts to proving that one has $K=Z$. Now this is true fiber by fiber by the hypothesis on $Z$, and it now suffices to apply IX~5.1 bis, which completes the proof of 4.3.

One will note that in criterion 4.3, the hypothesis that for every $s\in S$, $Z_s$ is the reductive center of $G_s$, is in fact purely geometric, \ie it suffices to verify it over the algebraic closure of $\kappa(s)$, as follows from the second assertion of 4.2.
\begin{theorem}{4.4}\label{XII.4.4}
Let $G$ be an affine algebraic group over a field $k$. Then $G$ admits a reductive center.
\end{theorem}
Since the center of $G$ is representable by a closed subgroup of $G$ (VIII~6.7), one is at once reduced to the case where $G$ is commutative. In addition, by 4.2 one can suppose $k$ algebraically closed. We shall see in XVII that in this case $G$ is written as a product $Z\times U$, where $Z$ is of multiplicative type and $U$ ``unipotent'', whence it will follow at once that $Z$ is indeed a reductive center of $G$.

We shall give here a proof independent of the structure theorem for commutative algebraic groups in the general form just indicated.

Note that it follows from IX~6.8 that the set of subgroups $H$ of multiplicative type of $G$ is increasingly filtered. Let us show that it admits a largest element.

When $G$ is smooth over $k$, one applies the structure theorem BIBLE~4 th.~4,
\[
G=G_s\times G_u,
\]
with $G_s$ of multiplicative type and $G_u$ ``unipotent'', which means here that $G_u$ admits a composition series whose factors are subgroups (smooth, if one wishes) of $\Ga$. (Indeed, $G_u/G_u^0$ is a unipotent group by BIBLE~4, cor. to th.~3, hence a $p$-group, where $p=\text{characteristic exponent}$, thanks to BIBLE~4 prop.~4; on the other hand $G_u^0$ admits a composition series with smooth connected quotients of dimension~1 by BIBLE~6 th.~1, cor.~1, and the latter are isomorphic to $\Ga$ by BIBLE~7 th.~4). Now one sees easily (\cf the lemma below) that every homomorphism from a group of multiplicative type to $\Ga$, hence also to $G_u$, is trivial, which does indeed prove that every subgroup of multiplicative type of $G$ is contained in $G_s$.

In the general case, consider the subgroup
\[
G_0=G_{\mathrm{red}}
\]
of $G$, which is smooth over $k$, hence by the preceding admits a largest subgroup $Z_0$ of multiplicative type. The subgroups of multiplicative type of $G$ containing $Z_0$ correspond to subgroups of multiplicative type of $G'=G/Z_0$. Now one has an exact sequence
\[
0\to G_0/Z_0\to G/Z_0\to G/G_0\to0,
\]
where, by the preceding, $G_0/Z_0$ has no subgroup of multiplicative type except the identity group. Since a subgroup of a group of multiplicative type is of multiplicative type (IX.8), it follows that for every subgroup $H$ of multiplicative type of $G/Z_0$, one has $H\cap(G_0/Z_0)=0$, hence $H\to G/G_0$ is injective. Since $G/G_0$ is a radicial algebraic group, hence finite over $k$, this implies that $H$ is itself radicial and of rank bounded above by that of $G/G_0$. This implies that the family of subgroups of multiplicative type of $G$ admits a largest element, say $Z$.

I say that $G/Z$ has no subgroup of multiplicative type other than the identity subgroup. This follows from the fact (proved in 7.1.1) that a commutative extension of two algebraic groups of multiplicative type is of multiplicative type: thus, if $Z'$ is the inverse image in $G$ of a subgroup of multiplicative type of $G/Z$, then $Z'$ is of multiplicative type, hence $Z'=Z$ by the maximal character of $Z$, so $Z'/Z=0$.

It now follows easily from the ``finite extension principle'' that for every extension $K$ of $k$, $(G/Z)_K=G_K/Z_K$ has no subgroup of multiplicative type except the identity group either.

We can now prove that $Z$ is a reductive center of $G$. Indeed, let $u:H\to G_S$ be a homomorphism of $S$-groups, where $S$ is a prescheme over $k$ and $H$ a group of multiplicative type and of finite type over $S$; let us prove that $u$ factors through $Z_S$. It amounts to the same thing to say that the composite homomorphism $H\to G_S\to(G/Z)_S=G_S/Z_S$ is zero.

Now I say that, putting $U=G/Z$, every homomorphism $u:H\to U_S$ is zero. Indeed, by IX~5.2 one is reduced to the case where $S$ is the spectrum of a field, and by IX~6.8 this then follows from the fact that $U_K$ has no subgroup of multiplicative type other than $0$. This completes the proof of 4.4. It remains only to set down the proof of the
\begin{lemma}{4.4.1}\label{XII.4.4.1}
Let $H$ be an $S$-group prescheme of multiplicative type; then every homomorphism of $S$-groups $u:H\to\Ga$ is trivial.
\end{lemma}
Indeed, consider the module $\OS^2=E$ as an extension of $\OS=E'$ by $\OS=E''$; then $\Ga$ is identified with the prescheme of automorphisms of this extension, so a homomorphism $u:H\to\Ga$ is identified with an $H$-module structure on $E$ respecting the extension structure, \ie such that $E'$ is stable under $H$ and that the actions induced by $H$ on $E'$ and $E''$ are trivial. By I~4.7.3 it follows that $H$ acts trivially on $E$, hence $u$ is trivial. \hfill Q.E.D.
\begin{remark}{4.4.2}\label{XII.4.4.2}
If $G$ is a nonzero abelian variety over $k$, then $G$ admits no reductive center in the sense of 4.1 with the restriction ``$G$ affine'' omitted, for, for $n$ prime to the characteristic, ${}_nG$ is étale over $k$ of order prime to the characteristic, hence of multiplicative type; but the family of ${}_nG$ is schematically dense in $G$, so if there were a reductive center, it would be identical to $G$, which is absurd, $G$ not being of multiplicative type. This is the reason for imposing in 4.1 the restriction that $G$ have affine fibers.
\end{remark}
\begin{lemma}{4.5}\label{XII.4.5}
Let $S$ be a prescheme, $G$ an $S$-group prescheme smooth over $S$, affine over $S$, with connected fibers, $T$ a maximal torus of $G$, $u:H\to G$ a central homomorphism, with $H$ an $S$-group prescheme of multiplicative type and of finite type. Then $u$ factors through $T$.
\end{lemma}
Indeed, let $C$ be the centralizer of $T$, which is a closed group sub-prescheme of $G$ smooth over $S$ (XI~5.3), hence affine over $S$. Since $T$ is in the center of $G$, it is invariant, and one can consider the quotient group $G/T=U$, which is representable (VIII~5.1). Since $u$ is central, it factors through $C$, and everything amounts to proving that the composite homomorphism $H\to C\to U=C/T$ is trivial. By IX~5.2 one is reduced to the case where $S$ is the spectrum of a field, which one can suppose algebraically closed. But then by BIBLE~6, th.~2, $U$ is a connected algebraic group (smooth over $k$) ``unipotent'', which means, as we have already observed, that it admits a composition series with quotients isomorphic to $\Ga$. Thus every homomorphism from a group $H$ of multiplicative type and of finite type to $U$ is trivial. This proves 4.5.
% typo?: "T is in the center of G" and G/T, later C/T, kept as printed.
\begin{corollary}{4.6}\label{XII.4.6}
Let $G$ be as in 4.5. If $G$ admits a reductive center, the latter is contained in every maximal torus of $G$.
\end{corollary}
\begin{theorem}{4.7}\label{XII.4.7}
Let $S$ be a prescheme, $G$ an $S$-group prescheme, smooth over $S$, affine over $S$, and with connected fibers.
\begin{enumerate}[label=\alph*)]
\item For every $s\in S$, let $z(s)$ be the type of the reductive center of $G_s$ (which is defined by 4.4). Order the set $E$ of isomorphism classes of $\ZZ$-modules of finite type by declaring that the class of $M$ is greater than that of $M'$ if $M'$ is isomorphic to a quotient of $M$. Then the map $S\to E$, $s\mto z(s)$ is lower semicontinuous.
\item For $G$ to admit a reductive center $Z$, it is necessary and sufficient that the preceding function $z:S\to E$ be locally constant. If this is so, $G/Z$ is representable (\cf VIII~5.1), and $G/Z$ admits the identity subgroup as reductive center.
\item Suppose that the reductive rank of $G$ is locally constant (\cf 1.7 b)). Then $G$ admits a reductive center $Z$. If $G/Z$ is representable (for example $G$ affine over $S$), then in addition the maximal tori $T$ of $G$ (resp. the Cartan subgroups $C$ of $G$) and $T'$ (resp. $C'$) of $G'=G/Z$ are in one-to-one correspondence, $T$ (resp. $C$) corresponding to $T'=T/Z$ (resp. $C'=C/Z$), and $T'$ (resp. $C'$) corresponding to $T=\varphi^{-1}(T')$ (resp. $C=\varphi^{-1}(C')$), where $\varphi:G\to G'$ is the canonical homomorphism.
\item Let $T$ be a maximal torus of $G$, denote by $\mathfrak g$ the Lie algebra of $G$, and consider the homomorphism
\[
\theta:T\to\GL(\mathfrak g),
\]
induced by the adjoint representation of $G$ (II~4). Then the kernel of $\theta$ is a reductive center of $G$.
\end{enumerate}
\end{theorem}
\begin{proof}
a) and b). The proof of a) and of the first assertion in b) is essentially identical to that of 1.7 a) and b), and we dispense with reproducing the argument here. Let us only point out that in a) one must appeal to IX~5.6 (using the fact that $G$ has connected fibers). Let us prove the second assertion of b), \ie that if $Z$ is a reductive center of $G$, then $G'=G/Z$ admits the identity subgroup as reductive center. Note at once that by VIII~5.1 the quotient group $G/Z$ is indeed representable; it is affine over $S$, of finite presentation over $S$ (VIII~5.8) and even smooth over $S$ (for it suffices to see this for the fibers, $G'$ being flat of finite presentation over $S$, but for the fibers it was pointed out in VI$_{\mathrm B}$.9.2.(xii) that a quotient of a smooth algebraic group over a field $k$ is smooth); finally, since $G$ has connected fibers, so does $G'$. Thus $G'$ satisfies the same initial hypotheses as $G$. To see that the identity subgroup of $G'$ is a reductive center, one can restrict oneself if one wishes to the case where $S$ is the spectrum of a field (4.3). Let $Z'$ be a central subgroup of multiplicative type of $G'$; everything amounts to proving that $Z'$ is reduced to the identity group. Let $Z_1$ be the inverse image of $Z'$ in $G$, and consider the action of $Z_1$ on $G$ induced by inner automorphisms of $G$. Since $Z$ is central, $Z$ acts trivially, so $Z_1$ acts through an action of $Z'$ on $G$. Moreover, since $Z$ is in the center of $G$, $Z_1$, hence $Z'$, acts trivially on $Z$; in addition, since $Z'$ is in the center of $G'$, the action of $Z'$ on the quotient $G/Z=G'$ is also trivial. Since $Z$ is central, it follows at once that the action of $Z'$ on $G$ comes from a group homomorphism
\[
u:Z'\to\uHom_{S\text{-gr}}(G',Z),
\]
by
\[
r(z')\cdot g=g\cdot u(z')(\varphi(g)),
\]
where $r:Z'\to\underline{\Aut}_{S\text{-gr}}(G)$ is the representation considered of $Z'$ by automorphisms of $G$, $\varphi:G\to G'$ the canonical homomorphism. Now giving a group homomorphism $u$ as above amounts to giving a group homomorphism
\[
v:G'\to\uHom_{S\text{-gr}}(Z',Z),
\]
and, on the other hand, by X~5.8 the second member is representable and is an étale group over $S$, hence its identity section is an open immersion, so $\Ker v$ is an open subgroup of $G'$. Since $G'$ has connected fibers, it is equal to $G'$, hence $v$ is zero, so $u$ is zero, so $Z'$ acts trivially on $G$, hence so does $Z_1$, which is therefore \emph{central} in $G$. Thus $Z_1$ is a \emph{commutative} extension of a group $Z'$ of multiplicative type by a group $Z$ of multiplicative type, hence (since we are over a field) is a group of multiplicative type (7.1.1). Being central in $G$, it is contained in the reductive center $Z$, \ie $Z_1=Z$, whence $Z'=\text{identity group}$, which completes the proof of b).

d) Let $Z=\Ker\theta$, which is a subgroup of multiplicative type of $T$ (for example by IX~6.8). By 4.5, every central homomorphism $u:H\to G$, with $H$ of multiplicative type and of finite type, factors through $T$, hence through $Z$. Since the formation of $Z$ is compatible with every base change, it remains to prove that $Z$ is central, \ie that the centralizer $C$ of $Z$ is equal to $G$. Now, by XI~5.3, $C$ is a smooth closed subgroup of $G$; on the other hand, since $Z$ acts trivially on $\mathfrak g$, one sees that $\Lie(C)=\mathfrak g$. One concludes easily that $C=G$: indeed, the immersion $C\to G$ is étale, for it is so fiber by fiber (as an unramified homomorphism of smooth algebraic groups of the same dimension, VI$_{\mathrm B}$.1.3), and one can apply X~3.5. Thus $C\to G$ is an étale closed immersion, hence an open immersion (SGA~1 I~5.1), and since it is also a closed immersion and $G$ has connected fibers, it is an isomorphism.

c) By 1.7 b), $G$ admits locally for the fpqc topology a maximal torus, so by 4.2 and d) just proved, $G$ admits a reductive center $Z$. We saw in d) that every maximal torus $T$ of $G$ contains $Z$. One observes at once that $T'=T/Z$ is a torus; to prove that it is a maximal torus in $G'$, one is reduced by definition 1.3 to the case where $S$ is the spectrum of an algebraically closed field, and then the assertion is contained in BIBLE~7 th.~3a). Conversely, let $T'$ be a maximal torus of $G'$, and let $T=\varphi^{-1}(T')$; let us prove that $T$ is a maximal torus of $G$. Since the question is local for the fpqc topology, one can suppose that $G$ already admits a maximal torus, say $T_0$, so that by the preceding, $T'_0=T_0/Z_0$ is a maximal torus of $G'$. By conjugacy theorem 1.7 c), $T'$ and $T'_0$ are locally conjugate for the fpqc topology, hence (since $\varphi:G\to G'$ is covering for this topology, so every section of $G'$ lifts locally to a section of $G$) $T$ and $T_0$ are likewise locally conjugate. Since $T$ is a maximal torus, so is $T$. One proceeds analogously for Cartan subgroups. This completes the proof.
% typo?: T'_0=T_0/Z_0 and "Since T is a maximal torus, so is T" kept as printed.
\end{proof}

Let us give a useful translation of d) in the case where $T$ is diagonalizable, hence of the form $D_S(M)$, where $M$ is a free $\ZZ$-module of finite type. Then (I~4.7.3) the $T$-module $\mathfrak g$ decomposes as a direct sum of sub-$T$-modules $\mathfrak g_m$ ($m\in M$):
\[
\mathfrak g=\coprod_{m\in M}\mathfrak g_m,
\]
which are necessarily locally free.

Suppose that for every $m\in M$, the rank of $\mathfrak g_m$ is constant (which is the case in particular if $S$ is connected). Then the set $R$ of $m\in M$ such that $\mathfrak g_m\ne0$ (``roots'') is finite. This being said:
\begin{corollary}{4.8}\label{XII.4.8}
Under the preceding conditions and with the preceding notation, the reductive center of $G$ is the intersection of the kernels of the root characters $m\in R$ on $T$. One therefore has an isomorphism
\[
Z\simeq D_S(N),
\]
where $N$ is the quotient of $M$ by the subgroup generated by $R$.
\end{corollary}
\begin{corollary}{4.9}\label{XII.4.9}
Let $G$ be an algebraic group over an algebraically closed field $k$. Suppose $G$ smooth over $k$, connected affine, with reductive center reduced to the identity group, and that the Lie algebra $\mathfrak g$ of $G$ is nilpotent. Then $G$ is ``unipotent'', \ie admits a composition series with quotients isomorphic to $\Ga$.
\end{corollary}
By BIBLE~6, th.~4, cor.~3, it suffices to prove that a maximal torus $T$ of $G$ is reduced to the identity group, or, what amounts to the same thing, that the Lie algebra $\mathfrak t$ of $T$ is reduced to $0$. Now it follows from the fact that the $T$-module $\mathfrak g$ decomposes according to the characters $\alpha$ of $T$ (I~4.7.3), that for every $t\in\mathfrak t$, the operation $\operatorname{ad}_{\mathfrak g}(t)$ on $\mathfrak g$ is semisimple. Thus if $\mathfrak g$ is nilpotent, $\operatorname{ad}_{\mathfrak g}(t)$ is zero. But by 4.7 d), since the reductive center of $G$ is reduced to the neutral element, the homomorphism $T\to\GL(\mathfrak g)$ is a monomorphism, hence induces an injective map on the Lie algebras, which means (II~4.5.) that for $t\in\mathfrak t$, the relation $\operatorname{ad}_{\mathfrak g}(t)=0$ implies $t=0$. This proves that $\mathfrak t=0$ and completes the proof.
\begin{proposition}{4.10}\label{XII.4.10}
Let $G$ be a smooth connected affine algebraic group over an algebraically closed field $k$. Then the reductive center of $G$ is the intersection of the maximal tori of $G$.
\end{proposition}
Of course, this is the intersection in the schematic sense (or, what amounts to the same thing, in the sense of subfunctors of $G$), \ie the largest closed sub-prescheme of $G$ contained in the maximal tori of $G$. It follows from the noetherian character of $G$ that it is also the intersection of a suitable \emph{finite} set of maximal tori of $G$.

Let $Z$ be the intersection in question; $Z$ is a closed subgroup of a torus, hence is of multiplicative type, and, on the other hand, by 4.5 it contains the reductive center of $G$. To prove that it is equal to it, it remains to prove that it is central. Since $G$ is connected, it suffices by IX~5.5 to prove that $Z$ is \emph{invariant}. Now by construction $Z$ is invariant under the $\operatorname{int}(g)$, with $g\in G(k)$, so the normalizer $N$ of $Z$ (\cf VIII~6.7) is a closed subgroup of $G$ containing the rational points of $G$. Since $G$ is reduced, one has $N=G$, which completes the proof.
\begin{proposition}{4.11}\label{XII.4.11}
Let $S$ be a prescheme, $G$ an $S$-group prescheme smooth over $S$, affine over $S$, with connected fibers, and of zero unipotent rank, which implies (1.9) that $G$ has locally constant reductive rank, hence (1.10) that the ``prescheme of maximal tori of $G$'' is defined, say $\mathscr T$, and is smooth, separated, of finite type over $S$.
Make $G$ act on $\mathscr T$ via inner automorphisms, whence a homomorphism of group functors
\[
u:G\to\underline{\Aut}_S(\mathscr T).
\]
Under these conditions, the following three subfunctors of $G$ are identical:
\begin{enumeratei}
\item The reductive center $Z$ of $G$.
\item The center $Z'$ of $G$.
\item The kernel $Z''=\Ker u$ of the preceding homomorphism.
\end{enumeratei}
In particular, the center of $G$ is representable by a subgroup of multiplicative type of $G$.
\end{proposition}

\begin{proof}
One trivially has $Z\subset Z'\subset Z''$; it remains to prove that $Z''\subset Z$, which amounts to proving (the hypotheses being stable under base change) that every section $g$ of $G$ over $S$ that acts trivially on $\mathscr T$ is a section of $Z$. Introducing $G'=G/Z$ and using 4.7 b) and c), which imply in particular that the prescheme $\mathscr T'$ of maximal tori of $G'$ is canonically isomorphic to $\mathscr T$, one is reduced to the case where $G=G'$, \ie to the case where the reductive center $Z$ of $G$ is zero. (N.B. Note that by 4.7 c), the unipotent rank of $G'$ is equal to that of $G$, hence is zero since that of $G$ is). Thus in this case one must prove that $g$ is the \emph{identity section} of $G$. The usual reduction procedure reduces us to the case where $S$ is noetherian, and even to the case where $S$ is \emph{artinian local} (since to verify that $g$ is the identity section, it suffices to verify, when $S$ is locally noetherian, that this is so after every base change $S'\to S$, with $S'$ artinian local). Now in this case the kernel $Z''$ of $u$ is representable (VIII~6.2 a) and 6.5 c)) (Note that $Z''$ is representable by XI~6.8). To prove that $Z''$ is reduced to the identity subgroup, it suffices by Nakayama to prove that this is so for its fiber $Z''_0$. This therefore reduces us to the case where $S$ is the spectrum of a field $k$, which one can obviously suppose algebraically closed. Now $Z''$ is contained in the stabilizer of every point of $\mathscr T(k)$, \ie in the normalizer $N(T)$ of every maximal torus $T$ of $G$. Since the reductive rank of $G$ is zero, it follows by XI~5.9 that $T$ is an open subgroup of $N(T)$, hence that the Lie algebra of $N(T)$ is identical to that of $T$. Thus the Lie algebra of $Z''$ is contained in that of $T$. On the other hand, it follows from 4.10 that the intersection of the Lie algebras of the maximal tori $T$ of $G$ is nothing other than the Lie algebra of the reductive center $Z$, hence here zero, since we supposed $Z=0$. Consequently the Lie algebra of $Z''$ is zero, \ie $Z''$ is étale over $k$. Moreover, $Z''$ is obviously invariant in $G$, and since $G$ is connected, it follows easily that $Z''$ is in the center of $G$. It is therefore in $T=\Centr_G(T)$ for every maximal torus $T$, hence in the intersection of the maximal tori, \ie in $Z=0$ by 4.10, which completes the proof.
% typo?: reductive rank zero, rather than unipotent rank zero, kept as printed.
\end{proof}

\par\addvspace{1.1\baselineskip}
\noindent{\large\bfseries 5.\ Application to the scheme of subgroups of multiplicative type{\renewcommand{\thefootnote}{*}\footnote{For a generalization of the results of the present \No to the case where $G$ is not supposed affine over $S$, \cf M.~Raynaud: \textit{Faisceaux amples sur les schémas en groupes et les espaces homogènes}, Chap.~IX.}\addtocounter{footnote}{-1}}}\par
\addvspace{0.45\baselineskip}
\addcontentsline{toc}{section}{5. Application to the scheme of subgroups of multiplicative type}
\begin{theorem}{5.1}\label{XII.5.1}
Let $S$ be a prescheme, $G$ an $S$-group prescheme smooth and affine over $S$, $M$ the ``scheme of subgroups of multiplicative type of $G$'', representing the functor described explicitly in XI~4.1. For every integer $n>0$, let $T_n$ be the subfunctor of $M$ such that $T_n(S')=$ set of group sub-preschemes $H$ of multiplicative type of $G_{S'}$ such that $n\mathrm{id}_H=0$, so $T_n$ is representable and is affine over $S$ (XI~3.12 a)). Let $u_n:M\to T_n$ be the morphism defined by $u_n(H)={}_nH$ (where ${}_nH=\Ker(n\mathrm{id}_H)$). With this notation, one has the following:
\begin{enumerate}[label=\alph*)]
\item Every sub-prescheme $U$ of $M$ of finite type over $S$ is contained in a closed sub-prescheme of finite type over $S$, and every closed sub-prescheme $X$ of $M$ of finite type over $S$ is affine over $S$.
\item Suppose $S$ quasi-compact, and let $X$ be a closed sub-prescheme of $M$ of finite type over $S$. Then there exists an integer $n>0$ such that for every multiple $m$ of $n$, the induced morphism
\[
u_m|_X:X\to T_m
\]
is a closed immersion.
\end{enumerate}
\end{theorem}
\begin{proof}
a) To prove the first assertion of a), one will take for $X$ the schematic closure of $U$ (EGA~I~9.5.1 and 9.5.3), which is defined since the immersion $i:U\to M$ is quasi-compact (for $U$ is of finite type over $S$ and $M$ separated over $S$ (4.1)). Thus one must prove that such an $X$ is affine over $S$, which will prove at the same time the second assertion of a). In the preceding form, one sees that the question is local on $S$, which one can therefore suppose affine. Then, since $U$ is of finite type over $S$, it is contained in a quasi-compact open subset of $M$, and this reduces us to the case where $U$ is itself an open subset of finite type over $S$, \ie quasi-compact. (N.B. A closed subscheme of a scheme affine over $S$ is affine over $S$).

It suffices to prove that such a $U$ is contained in a closed sub-prescheme of $M$ that is affine. In this form, the usual reduction procedure reduces us at once to the case where $S$ is noetherian. One proceeds likewise for b), which reduces to the case where $S$ is affine noetherian.

Take up again the projective limit $T$ of the $T_n$ used in the proof of XI~4.1, which is an affine prescheme (but not of finite type) over $S$, and whose local rings are noetherian rings, as was seen at the beginning of the proof of IX~3.7. (N.B. For $t\in T$, $t=(t_n)$, since the transition morphisms $T_m\to T_n$ are smooth and the dimension of the $T_m$ is bounded above, there exists an $n_0$ such that the $T_n\to T_{n_0}$ are étale at $t_n$ for every $n$ a multiple of $n_0$). The proof in \loccit, or XI~3.11, shows that the canonical morphism
\[
u:M\to T
\]
is an \emph{immersion}, and induces isomorphisms on the local rings (but one should take care that $u$ is not in general an open immersion or a quasi-compact morphism). Let $U$ be a quasi-compact open subset of $M$; we shall prove that its closure in $T$ is contained in $M$, which proves (if $X$ denotes the schematic closure of $U$ in $M$) that $u|_X:X\to T$ is a closed immersion, hence that $X$ is affine, and this will prove a).
% typo?: reference IX 3.7 kept as printed.

Since $U$ is noetherian, it has only finitely many irreducible components, and every element $t\in T$ in the closure of $U$ is a specialization of an element $x$ of $U$. Since the local ring of $t$ in $T$ is noetherian, so is the local ring $A$ of $t$ in $\overline x$ endowed with the induced reduced structure.

The canonical morphism from the integral noetherian local scheme $S'=\Spec(A)$ to $T$ then sends the closed point of $S'$ to $t$, the generic point to $x\in M$, and under these conditions one must show that $t\in M$. On replacing $A$ by the quotient of a suitable complete local $A$-algebra flat over $A$ (EGA~$0_{\mathrm{III}}$~10.3.1) by a minimal prime ideal, one can suppose that $A$ is complete with algebraically closed residue field (and one could even reduce to the case where it is in addition a discrete valuation ring, thanks to EGA~II~7.1.7). One is thus reduced (making the change of notation: $S'$ denoted by $S$) to the
\begin{lemma}{5.2}\label{XII.5.2}
Let $S$ be the spectrum of a complete integral noetherian local ring with algebraically closed residue field, $G$ an $S$-group prescheme smooth and affine over $S$, $(H(n))_{n>0}$ a system of subgroups of multiplicative type of $G$, indexed by the integers $n>0$, $\eta$ the generic point of $S$. Suppose:
\begin{enumerate}[label=\alph*)]
\item If $m$ is a multiple of $n$, one has $H(n)={}_nH(m)$.
\item There exists a subgroup $H_\eta$ of multiplicative type of $G_\eta$ such that one has $H(n)_\eta={}_n(H_\eta)$ for every $n>0$.
\end{enumerate}
Under these conditions, there exists a subgroup $H$ of multiplicative type of $G$ such that, for every $n>0$, one has $H(n)={}_nH$.
% typo: "te que" -> "tel que".
\end{lemma}
For every $n$, let $C_n=\Centr_G(H(n))$, which is a closed group sub-prescheme of $G$, smooth over $S$ (XI~5.3). The set of integers $n>0$ being ordered by divisibility, the $C_n$ form a decreasing family of closed sub-preschemes of $G$, and since $G$ is noetherian, this family is stationary. Let $C$ be the common value of the $C_n$ for large $n$. Then $C$ satisfies the same conditions as $G$; in addition the $H(n)$ are in fact subgroups of $C$. Thus, on replacing $G$ by $C$, one can suppose that the $H(n)$ are \emph{central} subgroups of $G$.

Let then $s$ be the closed point of $S$, and $Z_s$ the reductive center of $G_s$ (defined thanks to 4.4). By the definition (4.1), $Z_s$ contains the $H(n)_s$. Since $A$ is complete, $Z_s$ comes from a group sub-prescheme $Z$ of multiplicative type of $G$ (XI~5.8). I say that $Z$ contains the $H(n)$. Indeed, since $H(n)$ is central, hence commutes with $Z$ in $G$, one deduces a group homomorphism $Z\times H(n)\to G$, which, by IX~6.8, admits an image group that is a subgroup $K$ of multiplicative type of $G$ containing $Z$, and everything amounts to proving that $K=Z$, which follows from $K_s=Z_s$ and IX~5.1 bis.

One is thus reduced to proving the analogue of 5.2, but with $G$ replaced by a group $Z$ of multiplicative type and of finite type over $S$ (not necessarily smooth over $S$). Since $A$ is complete with algebraically closed residue field, $Z$ is diagonalizable (X~3.3 and 1.4), hence of the form $D_S(M)$, $M$ a commutative group of finite type. Consequently every subgroup $H$ of multiplicative type of $Z$ is diagonalizable (IX~2.11 (i)), hence of the form $D_S(N)$, where $N$ is a quotient group of $M$ (VIII~3.2). Thus the $H(n)$ correspond to quotient groups $M(n)$ of $M$, and the existence of a subgroup $H$ of multiplicative type of $Z$ such that $H(n)={}_nH$ for every $n$ amounts to that of a quotient group $N$ of $M$ such that $M(n)=N/nN$ for every $n$.

Now this follows at once from the fact that there exists a subgroup $H_\eta$ of $G_\eta$ such that $H(n)_\eta={}_n(H_\eta)$ for every $n$. This completes the proof of 5.2, hence of a).

b) We already know that $u|_X:X\to T$ is a closed immersion. It follows easily that for large $m$, the composite $u_m|_X:X\to T_m$ is also a closed immersion. Since we shall not need this fact later, we dispense with detailing its proof here.
\end{proof}
\begin{corollary}{5.3}\label{XII.5.3}
With the notation of 5.1, let $U$ be a subset both open and closed of $M$, of finite type over $S$. Then $U$ is affine over $S$ for the structure induced by $M$, and if $S$ is quasi-compact, there exists an integer $n>0$ such that for every multiple $m$ of $n$, the induced morphism $u_m|_U:U\to T_m$ is an open and closed immersion (\ie an isomorphism onto an open and closed subset of $T_m$, endowed with the induced structure).
\end{corollary}
The first assertion follows from 5.1 a), the second from 5.1 b), taking into account that $u_m:M\to T_m$ is smooth (XI~2.2 bis) and that a smooth, \ie étale, immersion is an open immersion.
\begin{corollary}{5.4}\label{XII.5.4}
Let $S$ be a prescheme, $G$ an $S$-group prescheme smooth and affine over $S$, of locally constant reductive rank (1.7 b)), $\mathscr T$ the ``prescheme of maximal tori of $G$'' (1.10). Then $\mathscr T$ is smooth and affine over $S$. If $T$ is a maximal torus of $G$, $N(T)$ its normalizer, then $G/N(T)$ (XI~5.3 bis) is affine over $S$. The same holds for $G/C(T)$ (where $C(T)$ is the centralizer of $T$), provided $W(T)=N(T)/C(T)$ is finite over $S$ (\cf 2.1 b)).
\end{corollary}
The second assertion is contained in the first, since by the conjugacy theorem, $G/N(T)$ is isomorphic to $\mathscr T$ (XI~5.5 bis). For the first assertion, one notes that by construction $\mathscr T$ is isomorphic to an open and closed sub-prescheme of $M$ (for one can suppose the reductive rank of $G$ constant and equal to $r$, and then $\mathscr T$ is the sub-prescheme of $M$ corresponding to subtori of relative dimension $r$, \ie the largest sub-prescheme of $M$ over which the ``universal subgroup of multiplicative type'' $H\subset G_M$ is a torus of relative dimension $r$). One can therefore apply 5.3. Finally, for the last assertion, one notes that $G/C(T)$ is finite over $G/N(T)\simeq(G/C(T))/W(T)$, hence is affine since $G/N(T)$ is.
\begin{corollary}{5.5}\label{XII.5.5}
Let $G$ be a smooth affine algebraic group over a field $k$. Then the scheme $M$ of subgroups of multiplicative type of $G$ is a direct sum of affine schemes over $S$. For every subgroup $H$ of multiplicative type of $G$, if $C$ and $N$ denote respectively its centralizer and its normalizer, the quotients $G/C$ and $G/N$ are affine.
% typo?: base S in a statement over k kept as printed.
\end{corollary}
Using XI~5.1 bis, one sees that the saturation under $G$ acting on $M$ of every finite closed subset of $M$ is open: indeed, one is reduced to the case where $k$ is algebraically closed, hence to the case of the orbit of a point $x$ rational over $k$, but then by \loccit the morphism $g\mto g\cdot x$ from $G$ to $M$ is smooth, hence open, so its image is open.

Let $U$ be the union of the classes of closed points of $M$ for the equivalence relation defined by the action of $G$. Then $U$ is open and contains every closed point of $M$, hence by the Nullstellensatz is identical to $M$. Thus $M$ is a disjoint union of open subsets, which are therefore necessarily closed, so $M$ is the sum prescheme of preschemes $M_i$, each of which is an orbit under $G$ of a closed point, hence is quasi-compact, hence of finite type. By 5.3 the $M_i$ are therefore affine. If $H$ is a subgroup of multiplicative type of $G$, it corresponds to a point $x$ of $M$ rational over $k$, and $G/N$ is identified with the orbit of $x$ under $G$ (XI~5.5 bis). It is therefore affine by the preceding. Since $C$ is an open subgroup of $N$ (XI~5.9), $G/C$ is therefore finite over $G/N$ (for the latter is isomorphic to $(G/C)/(N/C)$), hence affine since $G/C$ is. This proof shows at the same time:
% typo?: final G/C instead of G/N kept as printed.
\begin{corollary}{5.6}\label{XII.5.6}
Under the conditions of 5.5 for $G$ and $H$, the subscheme $U$ of $M$ ``of subgroups of multiplicative type of $G$ that are locally conjugate to $H$'' is an open and closed sub-prescheme of $G$.
% typo?: sub-prescheme of G instead of M kept as printed.
\end{corollary}
In pictorial terms, every subgroup $H'$ of multiplicative type of $G$ that is a limit of subgroups of $G$ conjugate to $H$ is itself conjugate to $H$.
\begin{remarks}{5.7}\label{XII.5.7}
Let $S$ be a prescheme, $G$ an $S$-prescheme smooth and affine over $S$, $H$ an $S$-prescheme of multiplicative type and of finite type; put $M=\uHom_{S\text{-gr}}(H,G)$, which by XI~4.2 is representable and is smooth over $S$ and separated over $S$. One can then prove for $M$ a result entirely analogous to 5.1, either by reducing to 5.1 by a proof analogous to that of XI~4.2, or by proceeding directly by a proof modeled on that of 5.1. One deduces corresponding variants for 5.3, 5.5, 5.6, which the reader will formulate.
\end{remarks}

\sgasection{6}{Maximal tori and Cartan subgroups of algebraic groups not necessarily affine (algebraically closed base field)}
\begin{lemma}{6.1}\label{XII.6.1}
Let $G$ be a connected algebraic group over a field $k$, $Z$ an algebraic subgroup of finite index of its center; then $G/Z$ is affine.
\end{lemma}
One can suppose that $Z$ is the center of $G$, for a scheme finite over an affine scheme is affine. Consider the vector spaces
\[
P^n=P^n(G)=\OO_{G,e}/\mathfrak m_{G,e}^{n+1}\qquad(n\text{ an integer }\ge0),
\]
where $\mathfrak m_{G,e}$ is the maximal ideal; then $G$ acts on the $P^n(G)$ by the adjoint representation, and if $Z_n$ is the kernel of the corresponding homomorphism
\[
\operatorname{ad}_n:G\to\GL(P^n),
\]
one verifies easily (using the fact that $G$ is connected) that $Z$ is the intersection of the $Z_n$, hence (since $G$ is noetherian) equal to one of the $Z_n$. But $\operatorname{ad}_n$ defines, by passage to the quotient, a monomorphism $G/Z_n\to\GL(P^n)$, which is therefore a closed immersion, hence $G/Z_n$ is affine, and consequently so is $G/Z$.
\begin{lemma}{6.2}\label{XII.6.2}
Let $G$ be a smooth algebraic group over the field $k$, $Z$ a central algebraic subgroup, $G'=G/Z$, $u:G\to G'$ the canonical homomorphism, $T$ a subgroup of multiplicative type in $G$, $T'=u(T)$ the image group, $C$ the centralizer of $T$ in $G$, $C'$ that of $T'$ in $G'$. Then one has
\[
C'{}^0\subset u(C).
\]
\end{lemma}
One can suppose $k$ algebraically closed. Let
\[
C_1=(u^{-1}(C')^0)_{\mathrm{red}};
\]
it suffices to prove that one has $C_1\subset C$, for $u(C_1)$ is connected of finite index in $C'$, hence equal set-theoretically to $C'{}^0$, hence equal to $C'{}^0$ since $C'$ and consequently $C'{}^0$ are smooth. Consider the morphism
\[
(g,t)\mto gtg^{-1}t^{-1}
\]
from $G\times G$ to $G$; it induces a morphism
\[
\varphi:C_1\times T\to Z_1,\quad\text{where }Z_1=Z_{\mathrm{red}},
\]
for, since the first member is reduced, it suffices to see that for $g\in C_1(k)$, $t\in T(k)$, one has $gtg^{-1}t^{-1}\in Z(k)$, which comes from the fact that $g$ centralizes $T$ modulo $Z$. One sees easily (by calculating on points rational over $k$) that $\varphi(g,t)$ is additive in $g$, and additive in $t$, hence is ``bilinear''; I say that (since $Z_1$ is smooth and $C_1$ connected) this homomorphism is identically zero, which will indeed prove that $C_1\subset C$. Using the density theorem for $T$, one is reduced to the case where $T$ is finite, \ie where there exists an integer $n>0$ such that $n\mathrm{id}_T=0$. Note that $\varphi$ is defined by a group homomorphism
\[
C_1\to\uHom_{S\text{-gr}}(T,Z_1),
\]
but, since $Z_1$ is commutative and smooth, the second member is representable by an étale algebraic group over $k$, and since $C_1$ is connected, every group homomorphism from $C_1$ to the latter is zero. \hfill Q.E.D.
% typo?: base S in Hom kept as printed.
\begin{corollary}{6.3}\label{XII.6.3}
Under the preceding conditions, suppose $C'$ connected; then one has
\[
C'=u(C),\qquad C=u^{-1}(C').
\]
\end{corollary}
Indeed, then $C'=C'{}^0$; on the other hand $C$ obviously contains $Z$, hence is equal to $u^{-1}(u(C))$.
\begin{lemma}{6.4}\label{XII.6.4}
Let $G$ be an algebraic group over the field $k$, $Z$ a central algebraic subgroup such that $G/Z=G'$ is a torus. Then $G$ is commutative, and if $k$ is algebraically closed, there exists a torus $T$ in $G$ such that $u(T)=G'$, where $u$ is the canonical homomorphism $G\to G/Z=G'$.
\end{lemma}
One can suppose $k$ algebraically closed. Consider again the morphism $G\times G\to G$ defined by $(g,h)\mto ghg^{-1}h^{-1}$; then (since $Z$ is central) this morphism factors through a morphism $G'\times G'\to G$, and, since $G'$ is commutative, the latter takes its values in $Z$, and even in $Z_1=Z_{\mathrm{red}}$, since $G'\times G'$ is reduced. One sees as above that the morphism $\varphi:G'\times G'\to Z_{\mathrm{red}}$ thus obtained is bilinear, hence zero, which proves that $G$ is commutative. To find a torus $T$ of $G$ such that $u(T)=G'$, one can (on replacing $G$ by $G_{\mathrm{red}}^0$) suppose $G$ smooth and connected, and in addition (on replacing $Z$ by $Z_{\mathrm{red}}^0$) that $Z$ is smooth and connected.

By a well-known theorem of Chevalley\nde{Give references here.}, $Z$ is an extension of an abelian variety by an affine group, which reduces us at once to proving our assertion in each of the following two cases: $1^\circ$) $Z$ is affine, $2^\circ$) $Z$ is an abelian variety. In case $1^\circ$), $G$ is affine and the result is well known (BIBLE~7 th.~3 a)). Thus suppose that $Z$ is an abelian variety. Since every homomorphism from the additive group $\Ga$ to the torus $G'=R$ or to the abelian variety $Z$ is trivial, it follows that the same holds for every homomorphism from $\Ga$ to $G$, hence $G$ contains no subgroup isomorphic to $\Ga$. By the structure theorem of Chevalley already invoked, one has an exact sequence
\begin{equation}\tag{*}
0\to T\to G\to A\to0,
\end{equation}
where $A$ is an abelian variety, and $T$ a smooth connected affine group. Since the latter is commutative and contains no additive subgroup, it follows that $T$ is a torus (and it is obviously the unique maximal torus of $G$). Everything amounts to proving that every epimorphism
\[
u:G\to R,
\]
where $R$ is a torus, satisfies $u(T)=R$. Put
\[
\Hom_{\mathrm{gr}}(T,\Gm)=M,\qquad\Hom_{\mathrm{gr}}(R,\Gm)=P,
\]
(these are free $\ZZ$-modules of finite type recovering $T,R$ up to isomorphism by $T=D_k(M)$, $R=D_k(P)$), and let
\[
B=\Ext^1_{k\text{-gr}}(A,\Gm),
\]
($B$ is also the set of points rational over $k$ of the abelian variety dual to $A$). One obviously has
\begin{equation}\tag{xx}
\begin{aligned}
\Ext^1(A,T)&=\Hom_{\mathrm{gr}}(M,B),\\
\Ext^1(A,R)&=\Hom_{\mathrm{gr}}(P,B),
\end{aligned}
\end{equation}
in particular the extension $G$ of $A$ by $T$ is given by a homomorphism
\[
\theta:M\to B.
\]
On the other hand, exact sequence (*) gives the exact sequence
\[
0\to\Hom(A,R)\to\Hom(G,R)\to\Hom(T,R)\to\Ext^1(A,R),
\]
and moreover $\Hom(A,R)=0$, and $\Hom(T,R)\to\Ext^1(A,R)$ is identified with the homomorphism
\[
\Hom(P,M)\to\Hom(P,B)
\]
deduced from $\theta:M\to B$. Putting
\[
N=\Ker(\theta),
\]
one therefore finds a canonical bijection
\[
\Hom(G,R)\simeq\Hom(P,N)\simeq\Hom(S,R),\qquad\text{where }S=D_k(N),
\]
which one can immediately describe geometrically as follows:
\begin{lemma}{6.5}\label{XII.6.5}
Let $G$ be an extension of an abelian variety $A$ by a torus $T=D_k(M)$, defined by a homomorphism $\theta:M\to B=\Ext^1(A,\Gm)$ (algebraically closed base field $k$). Let $N=\Ker\theta$, $S=D_k(N)$ the corresponding torus, isomorphic to $T/U$ where $U=D_k(M/N)$; consider the extension $H=G/U$ of $A$ by $S$. This extension splits\nde{Is split!}, so one has a unique projection from $H$ to $S$, whence a unique homomorphism
\[
v:G\to S
\]
extending the canonical homomorphism $T\to S$. This being said, for every torus $R$ and every homomorphism $u:G\to R$, there exists a unique homomorphism $u':S\to R$ such that $u=u'v$. In particular, one has $\Im(u)=\Im(u')=u(T)$.
\end{lemma}
This shows in particular that if $u$ is an epimorphism, so is its restriction to $T$, which completes the proof of 6.4.
\begin{theorem}{6.6}\label{XII.6.6}
Let $G$ be a smooth connected algebraic group over an algebraically closed field $k$.
\begin{enumerate}[label=\alph*)]
\item The maximal tori of $G$ are conjugate.
\item Let $T$ be a torus of $G$; then its centralizer is smooth and connected.
\item The map $T\mto\Centr_G(T)$ establishes a one-to-one correspondence between the set of maximal tori $T$ of $G$ and the set of Cartan subgroups (\No1) of $G$. For an algebraic subgroup $C$ of $G$ to be a Cartan subgroup, it is necessary and sufficient that it be smooth, nilpotent and set-theoretically equal to its connected normalizer (and then it is even equal to its connected normalizer); one then has $C=\Centr_G(T)$, where $T$ is the unique maximal torus of $C$, and $\Norm_G(C)=\Norm_G(T)$.
\item Let $v:G\to H$ be an epimorphism of smooth connected algebraic groups; then the maximal tori (resp. the Cartan subgroups) of $H$ are the images of the maximal tori (resp. the Cartan subgroups) of $G$. If $T$ is a maximal torus of $G$, $C$ its centralizer, then $v(C)$ is the centralizer of $v(T)$.
\item Under the conditions of d), suppose that $\Ker v$ is a central subgroup of $G$; then $T\mto v(T)$ (resp. $C\mto v(C)$) establishes a one-to-one correspondence between the set of maximal tori (resp. of Cartan subgroups) of $G$ and the set of maximal tori (resp. of Cartan subgroups) of $H$. The Cartan subgroups of $G$ contain the center of $G$ and a fortiori $\Ker v$, and are the groups of the form $v^{-1}(C')$, where $C'$ is a Cartan subgroup of $H$.
\end{enumerate}
\end{theorem}
Let us also describe explicitly at once the following immediate consequence of c):
\begin{corollary}{6.7}\label{XII.6.7}
For $G$ to be nilpotent (\ie the group $G(k)$ nilpotent), it is necessary and sufficient that the tori in $G$ be central, or again that $G$ have only one maximal torus (and then the latter is the largest subtorus of $G$).
\end{corollary}
\begin{proof}[Proof of 6.6]
Let $Z$ be a central algebraic subgroup of $G$, let $G'=G/Z$, and $u:G\to G'$ the canonical homomorphism. Then $G'$ is a smooth connected group. If $T'$ is a subtorus of $G'$, it follows from 6.4 that $u^{-1}(T')$ is commutative, and that $T'$ is the image of a subtorus of $u^{-1}(T')$, hence of a subtorus of $G$.

Since $u^{-1}(T')$ is commutative, it obviously admits a largest subtorus $T$ (for the sum of two subtori gives a third containing both), and one therefore has $u(T)=T'$. From this it follows immediately that for every maximal torus $T$ of $G$, its image $T'=u(T)$ is a maximal torus of $G'$, and that $T\mto u(T)$ is a one-to-one correspondence between maximal tori of $G$ and maximal tori of $G'$.

We now take
\[
Z=\Centr(G)^0_{\mathrm{red}},
\]
then $G'$ is affine by 6.1. Since the maximal tori of $G'$ are then conjugate, so are those of $G$, which proves a). Moreover, for $G$ to be nilpotent, resp. to have only one maximal torus, it is necessary and sufficient that $G'$ satisfy the same condition; but since $G'$ is affine, the two conditions in question on $G'$ are equivalent (BIBLE~6 th.~4 cor.~2), so the same holds for the conditions in question for $G$. Moreover, if $G$ has only one maximal torus $T$, the latter is invariant, hence central, and since every torus in $G$ is contained in a maximal torus, it is central. Conversely, if every torus is central, so are the maximal tori, and by conjugacy theorem a) there is only one maximal torus. This proves 6.7.

Let $T$ be any torus of $G$, $T'=u(T)$; then $C'=\Centr_{G'}(T')$ is connected (BIBLE~6 th.~6 a)), so by 6.3 the centralizer $C$ of $T$ is equal to $u^{-1}(C)$, hence connected (since $Z$ is connected), which proves b). If $T$ is maximal, hence $T'$ maximal, then one knows that $C'$ is nilpotent, hence $C$ (which is a central extension of $C'$) is nilpotent. In addition, $T$ is a maximal torus of $C$, hence by 6.7 it is the unique maximal torus of $C$, so the map $T\mto\Centr_G(T)$ from the set of maximal tori of $G$ to the set of Cartan subgroups is bijective. Moreover one has
\[
\Centr_G(T)=C\subset\Norm_G(C)\subset\Norm_G(T)
\]
and since one knows that the centralizer of $T$ is of finite index in its normalizer (\cf XI~5.9, whose argument is valid without an affine hypothesis, using only the representability of the two functors in question, as was pointed out in XI~6.5), and that $C$ is smooth and connected, one concludes
\[
C=\Norm_G(C)^0.
\]
In addition, by the one-to-one correspondence between maximal tori and Cartan subgroups, one sees that $\Norm_G(C)$ and $\Norm_G(T)$ have the same points with values in $k$, and since the second is smooth, one has
\[
\Norm_G(T)=\Norm_G(C).
\]
To complete the establishment of c), it remains to prove that if $C$ is a smooth connected nilpotent subgroup of $G$ of finite index in its normalizer, then $C$ is a Cartan subgroup. Now since $Z$ is central, the normalizer of $C$ contains $Z$, and since $Z$ is smooth and connected, one concludes $Z\subset C$, whence $C=u^{-1}(C')$, where $C'=u(C)$. One then has
\[
\Norm_G(C)=u^{-1}(\Norm_{G'}(C')),
\]
which proves that $C'$ is connected nilpotent of finite index in its normalizer, hence by BIBLE~7 th.~1 is a Cartan subgroup of $G'$, whence at once that $C$ is a Cartan subgroup of $G$.
% typo?: u^{-1}(C) rather than u^{-1}(C') retained above.

Let us prove e): one is under the conditions of the beginning of the proof (putting $\Ker u=Z$, $G'=H$); we have already seen that $T\mto u(T)$ is a one-to-one correspondence between maximal tori of $G$ and maximal tori of $G'$. Taking account of the one-to-one correspondence between maximal tori and Cartan subgroups just proved, one deduces a one-to-one correspondence between Cartan subgroups $C$ of $G$ and Cartan subgroups $C'$ of $G'$, by associating to $C=\Centr_G(T)$ the group $C'=\Centr_{G'}(T')$, where $T'=u(T)$. Since $C'$ is connected by b), it follows by 6.3 that $C'=u(C)$, and $C=u^{-1}(C')$, which proves e).

It remains to prove d). Let $T$ be a maximal torus of $G$; let us prove that $v(T)$ is a maximal torus of $H$ (which, taking account of conjugacy theorem a), will imply that the maximal tori of $H$ are all of the form $v(T)$ as above). Thus let $R$ be a torus in $H$ containing $v(T)$, and let us prove $R=v(T)$. On replacing $H$ by $R$, $G$ by $v^{-1}(R)^0_{\mathrm{red}}$, one can suppose $R=H$, \ie that $H$ is a torus, and one is reduced to proving that then $v(T)=H$. Let again $Z=\Centr(G)^0_{\mathrm{red}}$, $G'=G/Z$, and $H'=H/v(Z)$:
\begin{equation}\tag{D}
\xymatrix@C=1.1em{e\ar[r]&Z\ar[r]\ar[d]^{v''}&G\ar[r]^u\ar[d]^v&G'\ar[r]\ar[d]^{v'}&e\\e\ar[r]&v(Z)\ar[r]&H\ar[r]_{u'}&H'\ar[r]&e.}
\end{equation}
We already know that $u(T)$ is a maximal torus of $G'$, and since $G'$ is affine, hence $v':G'\to H'$ an epimorphism of smooth connected affine groups, $v'(u(T))$ is a maximal torus of $H'$ (BIBLE~7 th.~3 a)), hence equal to $H'$, \ie $u'(v(T))=H'$, so to prove $v(T)=H$, it suffices to show that $v(T)\supset v(Z)$. Now $v(Z)$ is a smooth connected subgroup of $H$, hence a torus, and $v'':Z\to v(Z)$ is an epimorphism, so by 6.4 one has $v(Z)=v''(S)$, where $S$ is a torus of $Z$. But $S$, being central in $G$, is obviously contained in the maximal torus $T$, whence $v(Z)\subset v(T)$. This proves assertion d) in the case of maximal tori.

Taking account of the one-to-one correspondence between maximal tori and Cartan subgroups, it remains to prove that if $T$ is a maximal torus of $G$, $C$ its centralizer, then $v(C)$ is the centralizer of $v(T)$. For this, take up again diagram (D) above (where of course $H$ is no longer supposed to be a torus), let $T'=u(T)$, $C'=u(C)$; we already saw in e) that $C'$ is the centralizer of the maximal torus $T'$, hence (since $G'$ is affine, hence $H'$ affine) $v'(C')$ is the centralizer of the maximal torus $v'(T')$ of $H'$ (BIBLE~7 th.~3 a)), \ie $u'(v(C))$ is the centralizer of $u'(v(T))$; since $C$ contains $Z$, hence $v(C)$ contains $v(Z)$, $v(C)$ is therefore the inverse image under $u'$ of $u'(v(C))$, \ie of the centralizer of $u'(v(T))$, hence is the centralizer of $v(T)$, as follows from e) applied to $u':H\to H'$ and to the maximal torus $v(T)$ of $H$. This completes the proof of 6.6.
% typo: missing closing parenthesis after BIBLE 7 th. 3 a) -> supplied.
\end{proof}

\sgasection{7}{Application to smooth group preschemes not necessarily affine}
\begin{theorem}{7.1}\label{XII.7.1}
Let $S$ be a prescheme, $G$ an $S$-group prescheme smooth, separated and of finite type over $S$. Suppose that $G$ admits locally for the faithfully flat quasi-compact topology a maximal torus. Then:
\begin{enumerate}[label=\alph*)]
\item The map
\[
T\mto\Centr_G(T)
\]
induces a bijection from the set of maximal tori of $G$ to the set of Cartan subgroups of $G$. If $C$ corresponds to $T$, then $T$ is the unique maximal torus of $C$.
\item Let $T,T'$ be two maximal tori of $G$, $C,C'$ the corresponding Cartan subgroups; then one has
\[
\uTransp_G(T,T')=\uTransp_G(C,C')=\uTransp_G(T,C'),
\]
the first two terms are also identical to the strict transporters, and finally the functor in question is representable by a closed sub-prescheme of $G$ smooth over $S$. The tori $T,T'$ and the Cartan subgroups $C,C'$ are locally conjugate for the étale topology.
\item There exists locally for the étale topology a maximal torus of $G$ and a Cartan subgroup of $G$.
\item Suppose that every finite subset of a fiber $G_s$ of $G$ is contained in an affine open subset of $G$ (for example $G$ quasi-projective over $S$, or $S$ artinian); then the functor $\mathscr T:\Sch^\circ_{/S}\to\Ens$ defined in 1.10 (the functor of maximal tori of $G$), isomorphic to the functor $\mathscr T:\Sch^\circ_{/S}\to\Ens$ of Cartan subgroups of $G$, is representable by a prescheme smooth, separated, of finite type over $S$, which is quasi-projective over $S$ when $G$ is, and is affine over $S$ when $G$ is affine over $S$, or when $S$ is artinian.
% typo?: both functors denoted T in source, retained.
\item Let $u:G\to G'$ be a morphism of $S$-group preschemes, where $G'$ is smooth, separated, of finite type over $S$, and suppose that for every $s\in S$, one has $u_s(G_s)=G'_s$, \ie $u_s$ is faithfully flat (Exp.~VI\nde{Reference not located in VI$_{\mathrm B}$; see M.~Demazure and P.~Gabriel, \textit{Groupes algébriques}, I, Masson (1970), proposition II.5.1.(c).}). Then for every maximal torus $T$ of $G$, $u(T)$ is a maximal torus of $G'$; if $G'$ has connected fibers, then for every Cartan subgroup $C$ of $G$, $u(C)$ is a Cartan subgroup of $G'$, and if $C$ is the centralizer of $T$, $u(C)$ is the centralizer of $u(T)$. In either case, the induced morphism $T\to u(T)$, resp. $C\to u(C)$, is faithfully flat.
\item Under the conditions of e), suppose in addition that $\Ker u$ is a central subgroup of $G$. Then $T\mto u(T)$ is a bijective map from the set of maximal tori of $G$ to the set of maximal tori of $G'$, and if $G'$ has connected fibers, $C\mto u(C)$ is likewise a bijective map from the set of Cartan subgroups of $G$ to the set of Cartan subgroups of $G'$; if $C'=u(C)$, one has $C=u^{-1}(C')$.
% typo: "une applications" -> "une application".
\end{enumerate}
\end{theorem}
\begin{remarks}{7.2}\label{XII.7.2}
We shall see in XV that the conclusion of d) remains valid without the restrictive condition on finite subsets of the $G_s$. We shall also prove there the conclusions concerning Cartan subgroups alone contained in b) c) d) e), when one supposes only that $G$ admits locally for the faithfully flat quasi-compact topology a Cartan subgroup (but not necessarily a maximal torus). For this, we shall moreover have to use 7.1 in the case where $S$ is artinian. Moreover, the proof of c) and d) (in the case of non-artinian $S$) is considerably simplified by using the method of XV.
% typo: "par nécessairement" -> "pas nécessairement".
\end{remarks}
\begin{proof}[Proof of 7.1]
a) Proceeding as in 3.2, one sees that for every maximal torus $T$ of $G$, $C=\Centr_G(T)$ is indeed a Cartan subgroup of $G$, and $T$ is determined in terms of $C$ as the unique maximal torus of $C$, so it remains only to show that every Cartan subgroup $C$ of $G$ is defined by a maximal torus of $G$, or again that $C$ admits a central maximal torus. Since the question is again local for the fpqc topology, one can suppose that $G$ admits a maximal torus $T$. Then by XI~6.2, $\uTransp_G(T,C)$ is representable by a closed sub-prescheme $S'$ of $G$ smooth over $S$; moreover $S'\to S$ is surjective, as follows from 6.6 c). Thus on making the base change $S'\to S$, one can suppose that there exists a section $g$ of $G$ such that $\operatorname{ad}(g)\cdot T\subset C$, but then $\operatorname{ad}(g)\cdot T$ is a maximal torus of $C$, which is central fiber by fiber by 6.6 c), hence central by IX~5.6 b). \hfill Q.E.D.

b) Let $T,T'$ be two maximal tori of $G$, and $C,C'$ the corresponding Cartan subgroups of $G$. Then the following conditions are equivalent:
\[
\begin{gathered}
(1)\ T\subset T'\qquad(2)\ T=T'\qquad(3)\ T\subset C'\\
(4)\ C\subset C'\qquad(5)\ C=C'.
\end{gathered}
\]
This follows trivially from a). Using the same result after arbitrary base change, one concludes the identity stated in b) between various transporters and strict transporters. Moreover, it was already observed in a) that $\uTransp_G(T,C')$ is representable by a closed sub-prescheme of $G$, smooth over $S$, and that its structural morphism is surjective. Consequently, by Hensel there exists locally for the étale topology a section of this prescheme over $S$, so $T$ and $T'$ on the one hand, $C$ and $C'$ on the other, are locally conjugate for the étale topology, which proves b).

c) Suppose first that $S$ is artinian local. When $G$ admits a maximal torus $T$, then it follows from the conjugacy theorem proved in b) that the functor $\mathscr T$ of maximal tori of $G$ is representable by the homogeneous space $G/N$, where $N$ is the normalizer of $T$ in $G$, which is smooth by b). Moreover, as was observed in VI$_{\mathrm A}$.3.2.1, since $\mathscr T$ is a homogeneous space under $G$, every finite subset of $\mathscr T$ is contained in an affine open subset. In the general case, there exists a finite extension $k'$ of the residue field $k$ of $S$ such that $G_{k'}$ has a maximal torus; then $k'$ comes from $k$ by a finite flat base change $S'\to S$, and the maximal torus of $G_{k'}$ lifts to a maximal torus $G_{S'}$ (XI~2.1 bis), so the functor $\mathscr T_{S'}$ is representable by a prescheme over $S'$, smooth, separated and of finite type over $S'$, every finite subset of which is contained in an affine open subset. Thus the natural descent datum on $\mathscr T_{S'}$ is effective, hence $\mathscr T$ is representable, and by descent one sees that $\mathscr T$ is smooth over $S$, separated and of finite type over $S$. From smoothness it follows, thanks to Hensel, that $\mathscr T$ admits a section locally for the étale topology.
% typo?: "a maximal torus G_{S'}" rather than "of G_{S'}" kept as printed.

This proves c) and d) in the case of artinian $S$ (N.B. We shall prove below that in this case $\mathscr T$ is in fact affine over $S$).

Suppose now $S$ arbitrary. To prove c) and d), which are local assertions on $S$ for the Zariski topology, one can suppose that there exists a prime number $\ell$ prime to the residue characteristics of $S$, that $S$ is affine, and that the reductive rank of the fibers of $G$ (which is obviously locally constant, thanks to the hypothesis of local existence for fpqc of a maximal torus) is constant, say $r$. Let $S'\to S$ be a faithfully flat quasi-compact morphism, $S'$ affine, such that $G_{S'}$ admits a maximal torus $T'$. Let $C'$ be its centralizer; I say that there exists a suitable power $n$ of $\ell$ such that one also has $C=\Centr_G({}_nT)$. Indeed, to see this, one is reduced at once to the case of noetherian $S'$, where this was seen in XI~6.2. Fix $n$ thus; put
\[
M=(\ZZ/n\ZZ)^r,\qquad P=\uHom_{S\text{-gr}}(M_S,G),
\]
then $P$ is obviously representable as a closed sub-prescheme of finite presentation of $({}_nG)^r$, where
\[
{}_nG=\uHom_{S\text{-gr}}((\ZZ/n\ZZ)_S,G),
\]
which is also the kernel of the $n$th-power morphism in $G$, hence representable by a closed sub-prescheme of finite presentation of $G$. Moreover $P$ is smooth over $S$ by XI~2.1. Let $s\in S$; by what was seen above, there exists a finite separable extension $k''$ of $\kappa(s)$ such that there exists a maximal torus $T''$ in $G_{k''}$; in addition, since ${}_nT''$ is étale over $k''$ ($n$ being prime to the characteristic of $k''$), one can suppose (on replacing $k''$ by a finite separable extension) that ${}_nT''$ is isomorphic to $M_{k''}$. Let $S''\to S$ be an étale morphism such that there exists an $s''\in S''$ above $s\in S$ giving rise to the residue extension $\kappa(s'')\simeq k''$. One therefore has a section of $P_{S''}\otimes\kappa(s'')$ over $\Spec(\kappa(s''))$, so by Hensel, on replacing $(S'',s'')$ by $(S''',s''')$ étale over it, one can suppose that there exists a section of $P_{S''}$ over $S''$, \ie an element of $P(S'')$, extending the given section. In other words, one has a homomorphism $M_{S''}\to G_{S''}$ inducing an isomorphism $M_{\kappa(s'')}\simeq{}_nT''$. By IX~6.4 (which here reduces to a simple application of Nakayama's lemma), this homomorphism is a closed immersion over an open neighborhood of $s'$, which one can suppose equal to $S''$. Let $H''$ be its image, and consider its centralizer $C''$ in $G_{S''}$, which is a smooth closed sub-prescheme of $G$, by XI~6.2. Note moreover:
% typo?: unprimed C,T in C=Centr_G(nT), and s' instead of s'', kept as printed.
\begin{lemma}{7.3}\label{XII.7.3}
Under the preceding conditions for $\ell,n,r$, for every prescheme $S''$ over $S$ and every maximal torus $T''$ of $G_{S''}$, one has
\[
\Centr_{G_{S''}}(T'')=\Centr_{G_{S''}}({}_nT'').
\]
\end{lemma}
Indeed, by faithfully flat descent from $S'\times_S S''$, one is reduced to the case where $G_{S''}$ admits a maximal torus $T''_1$ for which the preceding relation is true. But since $T''$ is locally conjugate to $T''_1$ for the étale topology by b), it follows that the same relation is true for $T''$.

Apply the preceding result with $\Spec(\kappa(s''))$ instead of $S''$; one sees that the fiber $C''_{s''}$ is a Cartan subgroup of $G$. Note now that since $G$ is smooth over $S$, the union of the neutral components of the fibers $G_s$ is an open subset of $G$ (by a general result of EGA~IV on smooth morphisms\nde{Make this reference precise.}), obviously stable under the group law of $G$, hence is a subgroup of $G$ for the prescheme structure induced by $G$.

In addition, $G^0$ satisfies the preliminary hypotheses on $G$, and there is an obvious one-to-one correspondence between the maximal tori of $G$ and those of $G^0$. Thus to prove c) and d), one can suppose $G$ with connected fibers, which we shall do. Then the Cartan subgroups of $G$ have connected fibers (6.6 a)). This being said, I say that $C''{}^0$ is a Cartan subgroup of $G_{S''}$ over an open neighborhood of $s''$ in $S''$ (which will complete the proof of c)). This follows indeed from the
\begin{lemma}{7.4}\label{XII.7.4}
Under the conditions of 7.1, suppose $G$ with connected fibers; let $C$ be a closed group sub-prescheme of $G$, smooth over $S$, and $s$ an element of $S$ such that $C_s$ is a Cartan subgroup of $G_s$; then $C^0$ is a Cartan subgroup of $G$ over an open neighborhood of $s$.
\end{lemma}
One reduces easily to the case where $S$ is local and $s$ its closed point, and to proving that $C$ is then a Cartan subgroup of $G$, then by flat descent to the case where $G$ admits a maximal torus, say $T$. Then by XI~6.2, $\uTransp_G(T,C)$ is representable by a closed sub-prescheme of $G$ smooth over $S$. In addition, by the hypothesis on $C_s$, the fiber of the transporter at $s$ is nonempty (taking account of conjugacy theorem 6.6 a)). This reduces us by faithfully flat descent to the case where this transporter admits a section over $S$, hence to the case where $C$ contains a maximal torus $T$ of $G$. But then $T$ is central in $C^0$ by IX~5.6 a), hence $C^0\subset\Centr_G(T)$, and since this is an inclusion of group schemes smooth over $S$, having the same relative dimension (namely the dimension of their common fiber at $s$) and connected fibers, it is an equality, which completes the proof.
% typo?: proving C rather than C^0 at the beginning kept as printed.

d) We retain the preceding notation and hypotheses for $\ell,n,r$, and the connectedness of the fibers of $G$. Let $Q:\Sch^\circ_{/S}\to\Ens$ be the functor defined by
\[
\begin{split}
Q(S')=\text{set of subgroups of multiplicative type of }G_{S'}\\
\text{of type equal to }(\ZZ/n\ZZ)^r=M\quad\text{(IX 1.4).}
\end{split}
\]
Then
\[
T\mto{}_nT
\]
is a morphism
\[
\varphi:\mathscr T\to Q,
\]
which is a monomorphism by 7.3. I say that $Q$ is representable by a separated prescheme of finite presentation over $S$. Indeed, as was pointed out in the proof of XI~3.12 a), one has an isomorphism
\[
Q\simeq P'/\Gamma
\]
where $P'$ is the open and closed sub-prescheme of the prescheme $P=\uHom_{S\text{-gr}}(M_S,G)$ introduced in c) corresponding to monomorphisms $M_{S'}\to G_{S'}$ (\cf IX~6.8), and where $\Gamma=(\Aut_{\mathrm{gr}}(M))_S$. The hypothesis that every finite subset of a fiber $G_s$ is contained in an affine open subset of $G$, being stable under passage to a closed sub-prescheme and under cartesian products, is obviously inherited by $({}_nG)^r$, hence by $P$, hence by $P'$, so $Q$ is representable by a prescheme of finite presentation over $S$ (\cf V)\nde{Make this reference precise.}.

One sees in the same way that if $G$ is quasi-projective (resp. affine) over $S$, so is $Q$. In any case, $Q$ is separated over $S$. Now one has the
\begin{lemma}{7.5}\label{XII.7.5}
The preceding homomorphism $\varphi=\mathscr T\to Q$ is representable by an open immersion.
% typo?: equals sign rather than colon kept as printed.
\end{lemma}
In other words, one must prove that if $H$ is a subgroup of multiplicative type of $G$, of type $M$, then there exists an open subset $U$ of $S$ such that for every $S'$ over $S$, $H_{S'}$ is of the form ${}_nT_{S'}$, for a suitable maximal torus $T_{S'}$ of $G_{S'}$, if and only if $S'\to S$ factors through $U$. One can obviously suppose that the nilpotent rank of the fibers of $G$ is constant (for, thanks to b), it is locally constant), say $r'$. Let $C=\Centr_G(H)$, which is a closed group sub-prescheme of $G$ smooth over $S$ (XI~6.2). Then, on replacing $S$ by the open and closed subset of points at which $C$ is of relative dimension $r'$, one can suppose $C$ of relative dimension $r'$ everywhere. One then sees at once that $H$ is of the form ${}_nT$, for a maximal torus $T$ of $G$, if and only if $C^0$ admits a central maximal torus $T$ of relative dimension $r$ everywhere and $H={}_nT$, which gives another expression for the subfunctor $U$ of $S$ that we want to represent (replacing $S$ in the preceding criterion by an $S'$ over $S$). Moreover, by flat descent one can suppose that $G$ admits a maximal torus $T_1$. Let $R$ be the subfunctor $\uTransp_G(T_1,\Centr_{C^0}(C^0))$ of $\uTransp_G(T_1,C)$; the latter is representable by a smooth prescheme over $S$ by XI~6.2, and the former is representable by an induced open sub-prescheme, as follows at once from IX~5.6 a); in particular it is smooth over $S$. Consequently its structural morphism to $S$ is open, hence its image is open, and on replacing $S$ by the said image (endowed with the induced structure), one can suppose the structural morphism surjective. Then, by 1.13 applied to $C^0$, $C^0$ admits a central maximal torus $T$, since it admits one locally for fpqc, which will obviously be of relative dimension equal to that of $T_1$, \ie $r$. Thus the condition to express for $H$ is the equality $H={}_nT$, which by IX~2.10 amounts again to taking a suitable open (and closed) subset of $S$.

Lemma 7.5 therefore implies that $\mathscr T$ is representable by a separated prescheme locally of finite presentation over $S$, and even of finite presentation over $S$, as one sees by taking up again the proof of 7.5 to ensure that $\varphi$ is in fact a quasi-compact open immersion, or by reducing by faithfully flat quasi-compact descent to the case where $G$ admits a maximal torus, and where $\mathscr T$ is therefore isomorphic to $G/\Norm_G(T)$. This last expression, or alternatively XI~2.1, shows in addition that $\mathscr T$ is smooth over $S$. Finally, if $G$ is quasi-projective over $S$, so is $Q$, hence also $\mathscr T$. If $G$ is affine over $S$, the assertion that $\mathscr T$ is then affine over $S$ is included as a reminder, having been established in 5.4 (N.B. I do not know whether, without an affine hypothesis for $G/S$, it is possible to choose $n$ so that in 7.5 the open immersion is also a closed immersion). For the assertion that $\mathscr T$ is affine over $S$ if $S$ is artinian, one is reduced to the case where $S$ is the spectrum of a field (EGA~I~6.1.7), which one can suppose algebraically closed. Then thanks to f), which will be proved below, it suffices to prove the same assertion for $G/\Centr(G)$; but the latter is affine by 6.1, so one is under the preceding conditions.

This completes the proof of d).

e) By IX~6.8, one knows that there exists a subtorus $T'$ of $G'$ such that $u$ induces a faithfully flat morphism $T\to T'$ (which characterizes $T'$ as the subsheaf $u(T)$ of $G'$). Let $C$ be the centralizer of $T$, $C'$ that of $T'$; let us prove that the morphism $C\to C'$ is flat, and faithfully flat if $G'$ has connected fibers. Since $C,C'$ are flat of finite presentation over $S$, one is reduced to the case of a base field (SGA~1 I~5.9), which one can obviously suppose algebraically closed. In addition one can suppose $G,G'$ connected (on replacing them by $G^0$ and $G'{}^0$, which does not change $C^0$ and $C'{}^0$), and it then suffices to apply 6.6 d), taking account of a).

f) Taking account of a) and e), one can restrict oneself to proving the assertion concerning Cartan subgroups. Now since a Cartan subgroup of $G$ is the centralizer of a maximal torus, it contains the center of $G$ and a fortiori $\Ker u$, so it is of the form $u^{-1}(C')$, where $C'=u(C)$ is the Cartan subgroup of $G'$ considered in e). Thus the map $C\mto u(C)$ is injective; to show that it is bijective, it suffices to see that for every Cartan subgroup $C'$ of $G'$, $u^{-1}(C')$ is a Cartan subgroup of $G$. Since the question is local for fpqc, one can suppose that $G$ admits a Cartan subgroup $C_1$, hence $u(C_1)=C'_1$ is a Cartan subgroup of $G'$, hence locally conjugate to $C'$ for fpqc by b), and since $u^{-1}(C'_1)=C_1$ is a Cartan subgroup of $G$, it follows that $u^{-1}(C')$ is a Cartan subgroup of $G$. \hfill Q.E.D.

One can also, to prove that $u^{-1}(C')$ is a Cartan subgroup, note that it is flat over $S$, for $u$ is (SGA~1 I~5.9), which reduces us by definition to the case of a base field, and one can apply 6.6 e).
\end{proof}
\begin{corollary}{7.6}\label{XII.7.6}
Under the conditions of 7.1 e), $G'$ admits locally for the étale topology a maximal torus, hence satisfies the preliminary conditions for $G$. If in addition $\Ker u$ is central, then the functors $\mathscr T_G,\mathscr C_G$ of maximal tori of $G$ and Cartan subgroups of $G$ are isomorphic to the analogous functors $\mathscr T_{G'},\mathscr C_{G'}$ for $G'$ (hence, in the case where they are representable, they are represented by isomorphic $S$-preschemes).
\end{corollary}
\begin{remark}{7.7}\label{XII.7.7}
a) Unlike what happens in the case where $G$ is affine over $S$ (it suffices, in fact, that $G$ have affine fibers, as one will see in XVI), it is not true that the fact that $G$ has locally constant reductive rank implies that $G$ admits locally for fpqc a maximal torus. Let, for example, $S$ be the spectrum of a discrete valuation ring, $G_1$ and $G_2$ smooth separated group schemes of finite type over $S$, such that the generic fiber of $G_1$ is an elliptic curve, the special fiber a group $\Gm$, and the generic fiber of $G_2$ a group $\Gm$, the special fiber a group $\Ga$, and take $G=G_1\times_S G_2$. Then the two fibers of $G$ have reductive rank~1, but one sees at once that $G$ admits no maximal torus locally for fpqc. On the other hand, it is very plausible that the following condition (for a smooth separated group of finite type over a prescheme $S$) is sufficient for the existence of a maximal torus locally for the étale topology: the reductive rank and the abelian rank of the fibers of $G$ are locally constant functions.

b) In the proof of 7.1 (in particular a)), we invoked XI~6.2 in cases where $S$ is not supposed locally noetherian. However, in the cases considered of application of XI~6.2, reduction to the case of affine noetherian $S$ is immediate.
\end{remark}
Here is a variant of 7.1 b):
\begin{proposition}{7.8}\label{XII.7.8}
Let $G$ be a smooth $S$-group prescheme of finite presentation with connected fibers, $H$ a group prescheme smooth over $S$, $i:H\to G$ a monomorphism of $S$-groups (making $H$ a sub-$S$-group of $G$). Then for every Cartan subgroup{\renewcommand{\thefootnote}{*}\footnote{The proof shows that it suffices to suppose that $C$ is a smooth subgroup of $G$ each of whose geometric fibers is the connected centralizer of a subgroup of a maximal torus of $G_{\overline S}$.}\addtocounter{footnote}{-1}} $C$ of $G$, $\uTransp_G(C,H)$ is representable by a closed sub-prescheme of $G$ smooth over $S$.
% typo?: capital S in G_{bar S} in the starred note kept as printed.
\end{proposition}
The fact that the transporter is representable by a closed sub-prescheme of $G$ of finite presentation is contained in XI~6.11, taking into account that $C$ has connected fibers since $G$ does (6.6 b)). To show that the transporter is smooth over $S$, one is reduced by the standard procedure to the case where $S$ is affine noetherian, then to the case where $S$ is artinian local, and by descent to the case where the residue field of $S$ is algebraically closed. But then $C$ admits a maximal torus $T$, which is a maximal torus of $G$. One can suppose that the reductive rank and the nilpotent rank of the fiber $H_0$ are equal to those of $G_0$ (otherwise the transporter would be empty), but then one sees at once (using the connectedness of $C$ and the fact that the centralizer in $H$ of a maximal torus of $H$ is smooth) that one has
\[
\uTransp_G(C,H)=\uTransp_G(T,H)
\]
and since one knows that the second member is smooth (XI~2.5), so is the first. The preceding argument proves more generally part b) of the
\begin{proposition}{7.9}\label{XII.7.9}
Let $G$ and $i:H\to G$ be as in 7.8; suppose in addition that for every $s\in S$, $H_s$ is connected, and has the same reductive rank and the same nilpotent rank as $G_s$ (\ie $H_s$ contains a Cartan subgroup of $G_s$). Then one has the following:
\begin{enumerate}[label=\alph*)]
\item $\uNorm_G(H)$ is representable by a closed sub-prescheme $\Norm_G(H)$ of $G$ of finite presentation over $S$, and the canonical monomorphism $H\to\Norm_G(H)$ is an open immersion; consequently $i$ is an immersion, and one has
\[
H=\Norm_G(H)^0.
\]
\item For every Cartan subgroup $C$ of $G$, $\uTransp_G(C,H)$ is a closed sub-prescheme of $G$, smooth over $S$, with surjective structural morphism. If $C$ is the centralizer of a maximal torus $T$ of $G$, one has in addition
\[
\uTransp_G(T,H)=\uTransp_G(C,H).
\]
\item Let $C$ be a group sub-prescheme of $H$. For it to be a Cartan subgroup of $H$, it is necessary and sufficient that it be a Cartan subgroup of $G$.
\item Suppose that $G$ admits locally for the étale topology, or for the fpqc topology, a Cartan subgroup (resp. a maximal torus); then the same holds for $H$.
\end{enumerate}
\end{proposition}
\begin{proof}
a) The representability of $\uNorm_G(H)$ by a closed sub-prescheme $\Norm_G(H)$ of $G$ of finite presentation over $S$ is contained in XI~6.11. Since $H$ is smooth, hence flat locally of finite presentation over $S$, to verify that $H\to\Norm_G(H)$ is an open immersion, one is reduced to verifying it on the fibers (VI$_{\mathrm B}$.2.6), which reduces us to the case where $S$ is the spectrum of an algebraically closed field $k$. One is then reduced (Exp.~VI\nde{Reference not located in VI$_{\mathrm B}$; see M.~Demazure and P.~Gabriel, \textit{Groupes algébriques}, I, Masson (1970), corollary II.5.6.}) to verifying that the corresponding homomorphism on the Lie algebras is an isomorphism, or, what amounts to the same thing, that $(\mathfrak g/\mathfrak h)^H=0$, where $\mathfrak g$ and $\mathfrak h$ are the Lie algebras of $G$ and $H$, and the exponent $H$ denotes the invariants under $H$ (\cf II~5.2.3 (i)). Now by hypothesis $H$ contains a Cartan subgroup $C$ of $G$, the centralizer of the maximal torus $T$ of $G$, and it therefore suffices to prove that one has
\[
(\mathfrak g/\mathfrak h)^T=0,
\]
which follows, taking account of the complete reducibility of representations of $T$ (I~4.7.3), from the analogous relation $(\mathfrak g/\mathfrak c)^T$, where $\mathfrak c=\Lie(C)$. As for the latter, equivalent to
\[
\mathfrak g^T=\mathfrak c,
\]
it means that the centralizer $C$ and the normalizer $N$ of $T$ have the same Lie algebra, which follows from the fact that $C$ is an open subgroup of $N$ (XI~5.9). This completes the proof of a).
% typo?: the relation (g/c)^T is printed without =0, retained.

b) As was pointed out, the proof was given in 7.8.

c) Taking account of the fact that $H$ is a sub-prescheme of $G$, the assertion reduces trivially to the case where $S$ is the spectrum of an algebraically closed field, in which case it follows at once from the hypothesis made on $H$.

d) One uses b), c) and ``Hensel's lemma'' XI~1.10.
\end{proof}
\begin{corollary}{7.10}\label{XII.7.10}
Let $G$ and $H$ be as in 7.9, and suppose that for every algebraically closed field $k$ over $S$, every element of $G_k(k)$ normalizing $H_k$ is in $H_k(k)$. Then $H$ is a closed sub-prescheme of $G$, and is its own normalizer.
\end{corollary}
This follows trivially from 7.9 a). We shall apply 7.10 in particular to the Borel subgroups (more generally, to the parabolic subgroups) of $G$.
\begin{corollary}{7.11}\label{XII.7.11}
Let $G,H$ be as in 7.9. Then $H$ contains every group sub-prescheme $Z$ of $G$, central in $G$ and flat and of finite presentation over $S$.
\end{corollary}
One can as usual reduce to the case of affine noetherian $S$, then to the case of artinian $S$, which implies that one is under the conditions of 7.1. By 7.9 c), one can suppose that $H$ contains a Cartan subgroup $C$ of $G$, and one is reduced to proving that $C\supset Z$. Now since $S$ is artinian, $G'=G/Z$ is representable by a group prescheme of finite type over $S$, the canonical morphism $u:G\to G'$ being faithfully flat and its kernel being $Z$ (VI$_{\mathrm A}$.3.2). Obviously $G'$ is smooth over $S$, and one can apply 7.1 f), which implies that $C$ is of the form $u^{-1}(C')$, hence contains $Z$.

\begin{corollary}{7.12}\label{XII.7.12}
Let $G,G'$ be two smooth $S$-group preschemes of finite presentation, $u:G\to G'$ a faithfully flat group homomorphism (\ie for every $s\in S$, $u_s$ is faithfully flat); suppose $\Ker u$ central and $G$ with connected fibers. Then the map
\[
H'\mto H=u^{-1}(H')
\]
establishes a one-to-one correspondence between group sub-preschemes $H'$ of $G'$, smooth of finite presentation over $S$, having the same reductive rank and the same nilpotent rank as $G'$ at every $s\in S$, and the set of group sub-preschemes $H$ of $G$, smooth of finite presentation over $S$, having the same reductive rank and the same nilpotent rank as $G$ at every $s\in S$. For $H$ to have connected fibers, it is necessary and sufficient that $H'$ do so.
\end{corollary}
Let $H$ be a group sub-prescheme of $G$ having the properties just specified. Then by 7.10, $H$ contains $Z=\Ker u$, hence by faithfully flat descent theory is of the form $u^{-1}(H')$, where $H'$ is a well-determined group sub-prescheme of $G'$, and one observes at once, taking account of 6.6 e), that the latter has the properties stated above. Moreover, if $H$ has connected fibers, so obviously does $H'=H/Z$. Thus it remains to prove that if one starts from a subgroup $H'$ of $G'$ having the stated properties, then $u^{-1}(H')=H$ has the same properties in $G$; and that if $H'$ has connected fibers, so does $H$. Taking account of 6.6 e), one is reduced to proving that $H$ is smooth over $S$ (resp. and has connected fibers). Now $H$ is already flat over $S$ as the inverse image of $H'$, which is so, under the flat morphism $u$, hence one is reduced to verifying that the geometric fibers of $H$ are smooth (resp. and connected), which reduces us to the case where $S$ is the spectrum of an algebraically closed field. Then $H'$ contains a Cartan subgroup $C'$ of $G'$, hence $H$ contains the inverse image $C$ of $C'$, which is a Cartan subgroup of $G$ by 6.6 e), hence $C$ is smooth and connected. Consequently 7.11 follows from the following lemma:
% typo?: references 7.10 and 7.11 in the proof of 7.12 kept as printed.
\begin{lemma}{7.13}\label{XII.7.13}
Let $G,G'$ be two flat group preschemes of finite presentation over $S$, $u:G\to G'$ a group homomorphism that is faithfully flat, $C'$ a group sub-prescheme of $G$ of finite presentation over $S$, such that $C=u^{-1}(C')$ is smooth over $S$ (resp. has connected fibers). Then for every group sub-prescheme $H'$ of $G'$ of finite presentation over $S$, containing $C'$, and such that $H'$ is smooth over $S$ (resp. has connected fibers), its inverse image $H=u^{-1}(H')$ is smooth over $S$ (resp. has connected fibers).
% typo?: C' of G rather than G' kept as printed.
\end{lemma}
As we observed above, this statement reduces at once to the case where $S$ is the spectrum of a field. Note then that $H$ is a principal bundle with base $H/C$ and group $C$ (Exp.~VI$_{\mathrm B}$.9); on the other hand $H/C\simeq H'/C'$ (\cf Exp.~IV), and since $H'$ is smooth (resp. connected), so is $H'/C'$, hence $H/C$. Since by hypothesis the same holds for $C$, it follows at once that it holds again for the bundle $H$. \hfill Q.E.D.

\sgasection{8}{Semisimple elements, union and intersection of maximal tori in group schemes not necessarily affine}
Throughout this \No, $G$ denotes a smooth group prescheme of finite presentation over $S$, with connected fibers.

Suppose first that $S$ is the spectrum of an algebraically closed field. When $G$ is affine, the notion of a \emph{semisimple element} of $G(k)$ was defined in BIBLE~4 \No4; one observes at once that this notion is invariant under an algebraically closed extension $k'/k$ of the base field. In addition, one saw in BIBLE~6 th.~5 (c) that $g\in G(k)$ is semisimple if and only if it is contained in a maximal torus of $G$. When $G$ is no longer supposed affine, $G$ is written canonically (thanks to Chevalley) as an extension of an abelian variety by a smooth connected affine algebraic group:
\begin{equation}\tag{*}
e\to V\to G\to A\to e.
\end{equation}
We shall say that an element $g$ of $G(k)$ is \emph{semisimple} if it is a semisimple element of $V(k)$. Since the maximal tori of $V$ are obviously identical to the maximal tori of $G$, it amounts to the same thing to say that $g$ belongs to a maximal torus of $G$. This is obviously again a notion invariant under an algebraically closed extension $k'/k$ of the base field $k$.

Suppose now that $S$ is the spectrum of an arbitrary field $k$, and let $g\in G$. Then (choosing an algebraically closed extension $K$ of $\kappa(g)$) one sees that $g$ is the image of a geometric point $g'$ of $G$ with values in an algebraically closed extension $K$ of $k$, and we shall say that $g$ is \emph{semisimple} if $g'$ is semisimple, which is independent of the particular choice of $g'$, thanks to what was said above. If $k'$ is an extension of $k$, then the set of semisimple elements of $G_{k'}$ is the inverse image of the set $G^{\mathrm{ss}}$ of semisimple elements of $G$.

Suppose finally $S$ arbitrary; then a point $g\in G$ is said to be semisimple if it is semisimple in its fiber $G_s$. If $S'\to S$ is any morphism, then $G^{\mathrm{ss}}_{S'}$ is therefore the inverse image of $G^{\mathrm{ss}}$. Suppose that the functor $\mathscr T$ defined in 1.10 (the functor of maximal tori) is representable by a prescheme of finite presentation over $S$ (which is the case, for example, if $G$ admits locally for fpqc a maximal torus, by 7.1 d), at least if $G$ is quasi-projective over $S$). Consider the canonical maximal torus $\underline T$ of $G_{\mathscr T}$ (``universal maximal torus of $G$''), and the morphism
\[
u:\underline T\to G
\]
induced by the projection $G_{\mathscr T}\to G$. Then it follows at once from the definition that $G^{\mathrm{ss}}$ is nothing other than the image of the preceding morphism. One will conclude:
\begin{proposition}{8.1}\label{XII.8.1}
The set $G^{\mathrm{ss}}$ of semisimple elements of $G$ is locally constructible (hence constructible if $S$, hence $G$, is quasi-compact and quasi-separated).
\end{proposition}
One reduces as usual to the case where $S$ is affine noetherian; in addition one can suppose (by the usual noetherian criterion of constructibility (EGA~$0_{\mathrm{III}}$~9.2.3)) that $S$ is integral, and restrict oneself to proving that there exists a nonempty open subset $U$ of $S$ such that $G^{\mathrm{ss}}|_U$ is constructible. Taking $U$ small enough, and replacing it if necessary by a finite covering, one can suppose that $G$ is separated over $S$ and contains a maximal torus $T$. But then by 7.1 d), the functor $\mathscr T$ is representable by a prescheme of finite presentation over $S$, and so therefore is $\underline T$, whose image in $G$ is consequently constructible. \hfill Q.E.D.

Suppose again that $k$ is a field, and consider the subgroup $H$ of $G$ generated by the preceding morphism (\cf VI$_{\mathrm B}$.1.2). It is a smooth group subscheme of $G$, connected since $\underline T$ is, whose formation is obviously compatible with every extension of the base field (\cf VI), that of $\underline T$ being so. When $k$ is algebraically closed, one sees immediately that $H$ is also the algebraic subgroup of $G$ generated by the maximal tori of $G$, or, what amounts to the same thing, by the tori of $G$, and it is also the smallest algebraic subgroup of $G$ containing the semisimple elements of $G(k)$. (In fact, these characterizations of $H$ remain valid as soon as $k$ is an infinite field, thanks to the fact, proved in XIV, that the set of points of $\underline T$ rational over $k$ is dense in $\underline T$). Moreover, $H$ is invariant in $G$, for to see this one can restrict oneself to the case where $k$ is algebraically closed, and then ($H$ and $G$ being smooth over $k$) it suffices to verify that $H$ is stable under the inner automorphisms $\operatorname{int}(g)$, $g\in G(k)$, which is obvious. (One could even show that $H$ is a characteristic subgroup of $G$, \ie stable under $\underline{\Aut}_{k\text{-gr}}(G)$). It then follows at once from 6.6 d) that the reductive rank of $G/H$ is zero; more precisely, if $K$ is an invariant algebraic subgroup of $G$, it follows at once from the fact that for algebraically closed $k$ the maximal tori of $G/K$ are the direct images of the maximal tori of $G$, that the reductive rank of $G/K$ is zero if and only if $K$ contains all the maximal tori of $G$ ($k$ being supposed algebraically closed), or again if and only if $K$ contains $H$ ($k$ arbitrary): hence $H$ is the smallest invariant algebraic subgroup of $G$ such that $G/H$ has zero reductive rank. Another obvious characterization of $H$ is the following: it is the smallest algebraic subgroup of $G$ having the same reductive rank as $G$. Note finally that $H$ is affine: indeed, to see this one can again suppose $k$ algebraically closed, and taking up again exact sequence (*) from the beginning of the \No, one notes that the maximal tori of $G$ are contained in $V$, hence so is $H$, so since $V$ is affine, $H$ is. Let us summarize the results obtained:
\begin{proposition}{8.2}\label{XII.8.2}
Let $G$ be a smooth connected algebraic group over the field $k$, and let $\overline k$ be an algebraic closure of $k$. There exists an algebraic subgroup $H$ of $G$ such that $H_{\overline k}$ is the algebraic subgroup of $G_{\overline k}$ generated by the maximal tori (or again by the tori) of $G_{\overline k}$. The group $H$ is also characterized as the smallest algebraic subgroup of $G$ having the same reductive rank as $G$, or the smallest invariant algebraic subgroup of $G$ such that the reductive rank of $G/H$ is zero. It is a smooth connected invariant affine subgroup of $G$, whose formation commutes with every extension of the base field.
\end{proposition}
To use the characterization of $H$ in terms of $G/H$, it is appropriate to state explicitly the
\begin{corollary}{8.3}\label{XII.8.3}
Let $G$ be a smooth connected algebraic group over an algebraically closed field $k$. For $G$ to have zero reductive rank (\ie with the notation of 8.2, for one to have $H=(e)$), it is necessary and sufficient that $G$ be an extension of an abelian variety by a smooth connected unipotent algebraic group (\ie a successive extension of groups isomorphic to the additive group $\Ga$).
% typo: duplicated k in "un corps k algébriquement clos k" -> single k.
\end{corollary}
Indeed, thanks to Chevalley's exact sequence (*), one is reduced to proving that the reductive rank of the smooth connected affine group $V$ is zero if and only if $V$ is unipotent, which is contained in BIBLE~6.4. th.~4 cor.~3.

Thus, with the notation of 8.2, $H$ is the smallest invariant algebraic subgroup of $G$ such that $(G/H)_{\overline k}$ is an extension of an abelian variety by a smooth connected affine unipotent algebraic group. One also concludes:
\begin{corollary}{8.4}\label{XII.8.4}
With the notation of 8.2, for one to have $G=H$ (\ie $G_{\overline k}$ generated by its maximal tori), it is necessary and sufficient that $G$ be affine and that every homomorphism from $G_{\overline k}$ to the additive group be trivial.
% typo: "la groupe additif" -> "le groupe additif".
\end{corollary}
\begin{remarks}{8.5}\label{XII.8.5}
a) Let $\overline V$ be the largest smooth connected affine algebraic subgroup of $G_{\overline k}$ (so $G_{\overline k}$ is an extension of an abelian variety by $\overline V$). It is well known (Rosenlicht\nde{Indicate references here.}), if $k$ is not perfect, that $\overline V$ is not in general ``defined over $k$'', \ie there does not in general exist an algebraic subgroup $V$ of $G$ such that $V_{\overline k}=\overline V$. However, when $\overline V$ is generated by its maximal tori, \ie when $\overline V$ has no quotient group isomorphic to the additive group $\mathbf G_{\mathrm a,\overline k}$, there exists such a $V$, namely the group $H$ of 8.2. One therefore sees that in this question of rationality, as in many others (\cf for example XIV \No6), all the troubles come from unipotent groups, \ie from the additive group, whereas it is the presence of (sufficiently many) groups of multiplicative type that ensures, on the contrary, that things work well.

b) One can also introduce, with the notation of 8.2, the inverse image $H'$ in $G$ of the commutator subgroup of $G/H$; then $H'$ is the smallest invariant algebraic subgroup of $G$ such that $(G/H')_{\overline k}$ is a commutative extension of an abelian variety by a smooth connected affine unipotent algebraic group. $H'$ is again a smooth connected affine subgroup of $G$. Let $H''$ be the inverse image in $G$ of $p^n(G/H')$ for large $n$, $p$ being the characteristic (supposed $>0$); then one sees easily that $H''/H'$ is an abelian variety, and $H''$ is the smallest invariant algebraic subgroup of $G$ such that $G/H''$ is a smooth connected affine unipotent commutative algebraic group. This being said, one sees easily that for $\overline V=H_{\overline k}$ (notation of a)), \ie for every homomorphism from $\overline V$ to the additive group to be trivial, it is necessary and sufficient that $H''=G$, or again that every homomorphism from $G_{\overline k}$ to the additive group be trivial.
\end{remarks}

To finish, we shall generalize to smooth groups with connected fibers the notion of reductive center developed in \No4, taking inspiration from 4.10. Let $Z$ be the subfunctor of $G$ defined by
\[
\begin{aligned}
Z(S')={}&\text{Set of sections }g'\text{ of }G_{S'}\text{ over }S'\text{ such that}\\
&\text{for every }S''\text{ over }S'\text{ and every maximal torus }T''\\
&\text{of }G_{S''},\text{ the inverse image }g''\text{ of }g'\text{ under }S''\to S'\\
&\text{is a section of }T''\text{ over }S''.
\end{aligned}
\]
Introducing the functors $\mathscr T(S')=$ set of maximal tori of $G_{S'}$, and $\underline T(S')=$ set of pairs $(T',g')$, where $T'$ is a maximal torus of $G_{S'}$ and $g'$ a section of $T'$ over $S'$, one sees that $\underline T$ is a subfunctor of $\mathscr T\times_S G=\mathscr T_G$, and with this notation one can also write
\[
Z=\prod_{\mathscr T_G/G}\underline T/\mathscr T_G.
\]

Using XI~6.8, and 7.1 d), which ensures the representability of $\mathscr T$ by a smooth $S$-prescheme under certain conditions, one could conclude the representability of $Z$ under certain conditions, which we shall, however, obtain by a more direct route below.
\begin{definition}{8.6}\label{XII.8.6}
Let $G$ be a smooth $S$-group prescheme of finite presentation with connected fibers. One says that $G$ \emph{admits a reductive center} if the preceding functor $Z$ (which is obviously a subgroup of $G$) is representable by a group of multiplicative type. One then says that $Z$ is the \emph{reductive center} of $G$.
\end{definition}
One will note that if $Z$ is a reductive center of $G$, then for every base change $S'\to S$, $Z_{S'}$ is a reductive center of $G_{S'}$; on the other hand, the existence of a reductive center is obviously a local question for the fpqc topology. As for the terminology ``reductive center'', note that $Z$ is in any case central, for obviously $Z$ is invariant under
\[
\underline{\Aut}_S(G)
\]
and a fortiori is an invariant subgroup of $G$, and one applies IX~5.5.

Lemma 4.5 must here be replaced by:
\begin{lemma}{8.7}\label{XII.8.7}
Let $u:H\to G$ be a central homomorphism, where $H$ is of multiplicative type and of finite type over $S$; suppose that for every algebraically closed field $k$ over $S$, $u_k(H_k)$ is contained in the largest smooth connected affine subgroup of $G_k$. Then $u$ factors through every maximal torus $T$ of $G$ (hence $u$ in fact factors through the subfunctor $Z$ of $G$ defined above).
\end{lemma}
Using an easy variant of IX~5.1 bis (where the sign $=$ would be replaced by an inclusion sign), one is reduced to the case where $S$ is the spectrum of a field $k$, which one can suppose algebraically closed. (Reduce to the case of affine noetherian $S$, then artinian, then use IX~3.6). Since $T$ is contained in the largest smooth connected affine subgroup $V$ of $G$, one is then reduced to the case where $G=V$, \ie where $G$ is affine, where the result was proved in 4.5.
\begin{proposition}{8.8}\label{XII.8.8}
Let $G$ be a smooth $S$-group prescheme of finite presentation with connected fibers.
\begin{enumerate}[label=\alph*)]
\item If $S$ is the spectrum of a field $k$, then $G$ admits a reductive center $Z$. When $G$ is an extension of an abelian variety by a smooth connected affine algebraic group $V$ (for example if $k$ is algebraically closed), then $Z$ is also the reductive center of $V$, and it is the largest central subgroup of multiplicative type of $V$.
\item Let $Z$ be a group sub-prescheme of multiplicative type of $G$. Then $Z$ is a reductive center of $G$ if and only if for every $s\in S$, $Z_s$ is a reductive center of $G_s$. Then $Z$ is the largest subgroup $K$ of multiplicative type of $G$ such that for every $s\in S$, $K_s$ is contained in the reductive center of $G_s$; more generally, for every homomorphism $u:H\to G$, with $H$ of multiplicative type and of finite type over $S$, such that for every algebraically closed field $k$ over $S$, $u_k:H_k\to G_k$ factors through the largest smooth connected affine subgroup of $G_k$, $u$ factors through $Z$ (and in particular $u$ is central).
% typo?: no centrality hypothesis on u in this assertion, retained.
\item If $G$ admits locally for the fpqc topology a maximal torus, then $G$ admits a reductive center.
\item Let $T$ be a maximal torus of $G$. Then $T\cap\Centr(G)=\Ker(T\to\GL(\mathfrak g))$ and this is a reductive center of $G$.
\item Let $Z$ be a reductive center of $G$, and suppose $G'=G/Z$ representable (for example $S$ artinian); then $G'$ admits the identity subgroup as reductive center, and $T'\mto T=u^{-1}(T')$ establishes a one-to-one correspondence between maximal tori of $G'$ and maximal tori of $G$.
\end{enumerate}
\end{proposition}
\begin{proof}
a) Suppose that $S$ is the spectrum of a field $k$. To prove then the existence of a reductive center, one can suppose $k$ algebraically closed, and consequently one is reduced to the case where $G$ is an extension of an abelian variety by a smooth connected affine algebraic group $V$. Since for every $S'$ over $k$, the maximal tori of $G_{S'}$ are those of $V_{S'}$ (by IX~5.2 and the fact that over an algebraically closed field a homomorphism from a torus to an abelian variety is trivial), it follows that the functor $Z$ defined above in terms of $G$ is the same as that defined in terms of $V$. One is thus reduced to the case of affine $G$. Since $\mathscr T$ is representable, and necessarily ``essentially free'' over $k$ (VIII~6.1), it follows that $\mathscr T_G$ is essentially free over $G$, and since $\underline T$ is a closed subscheme of $\mathscr T_G=G_{\mathscr T}$ (by, for example, VIII~5.7), it follows by VIII~6.4 that $Z=\prod_{\mathscr T_G/G}\underline T/\mathscr T_G$ is representable by a closed subscheme of $G$. Thus it is a group subscheme of $G$; I say that it is of multiplicative type: indeed, one can suppose $k$ algebraically closed; then $G$ admits a maximal torus $T$, and by definition one will have $Z\subset T$, so $Z$ is of multiplicative type as an algebraic subgroup of a group of multiplicative type (IX.8). This proves that $Z$ is a reductive center of $G$. The fact that it is the largest central subgroup of multiplicative type of $V$ is contained in 8.7.

b) The ``only if'' being trivial, let us prove that if $Z$ is a subgroup of multiplicative type of $G$ such that for every $s\in S$, $Z_s$ is the reductive center of $G_s$, then $Z$ is a reductive center of $G$. One must first prove that $Z$ is contained in every maximal torus of $G$ (which will therefore remain true after every base change): this is an immediate consequence of 8.7. Next, one must prove that if $g$ is a section of $G$ over $S$ such that for every $S'$ over $S$, and every maximal torus $T'$ of $G_{S'}$, $g_{S'}$ is a section of $T'$, then $g$ is a section of $Z$. Note that one could reduce as usual (taking account of the fact that $\prod_{\mathscr T_G/G}\underline T/\mathscr T_G=Z'$ is a sheaf for fpqc that commutes with inductive limits of rings) to the case of affine $S$, then noetherian $S$, and finally artinian $S$. Then $\mathscr T$ is representable by a smooth prescheme over $S$ by 7.1 d), so proceeding as in a), one sees that $\prod_{\mathscr T_G/G}\underline T/\mathscr T_G$ is representable by a closed subscheme $Z'$ of $G$. We have already seen that $Z\subset Z'$; on the other hand, by the hypothesis on $Z$, one has $Z_k=Z'_k$ (where $k$ is the residue field), but since $Z$ is flat over $S$, it follows that $Z=Z'$, which proves that $Z$ is a reductive center of $G$. The other assertions of b) are contained in 8.7.

c) Reduces immediately to d).

d) Of course, $\mathfrak g$ denotes the Lie algebra of $G$, and $T\to\GL(\mathfrak g)$ the homomorphism induced by the adjoint representation of $G$. One trivially has
\[
T\cap\Centr_G\subset\Ker(T\to\GL(\mathfrak g));
\]
the proof of the reverse inclusion is the same as in 4.7 d), and we do not repeat it here. Let $Z$ be the group in question; as $\Ker(T\to\GL(\mathfrak g))$ it is of multiplicative type (\cf for example IX~6.8). To prove that it is a reductive center of $G$, one is reduced by b) to the case where $S$ is the spectrum of a field. Let then $Z'$ be the reductive center (which exists by a)); one obviously has $Z'\subset T\cap\Centr(G)=Z$; on the other hand, since $Z$ is a central subgroup of multiplicative type contained in the smooth connected affine subgroup $T$, it follows from a) that $Z'\supset Z$, hence $Z'=Z$. \hfill Q.E.D.

e) Under the conditions of 7.1 f), put $Z=\Ker u$ and suppose that for every algebraically closed field $k$ over $S$, $Z_k$ is contained in a maximal torus of $G_k$. Then one observes easily that the map $T'\mto u^{-1}(T')$ induces a one-to-one correspondence between the set of maximal tori of $G'$ and the set of maximal tori of $G$. Applying this to situation 8.8 e), the desired conclusion follows at once.
\end{proof}
\begin{remarks}{8.9}\label{XII.8.9}
a) The proof given of 8.8 is independent of the results of \No4, and in particular that of 8.8 a) does not use 4.4 (whose proof is somewhat laborious).

b) One sees easily that the subgroup $Z$ of $G$ considered in 8.6 is always central (whether representable or not), and it may be tempting to call it the reductive center of $G$ in all cases.

c) One can also generalize 4.9; one finds the following statement: let $G$ be a smooth connected algebraic group over an algebraically closed field; for $G$ to be an extension of an abelian variety by a smooth connected unipotent algebraic group (\ie for the reductive rank of $G$ to be zero, \cf 8.3), it is necessary and sufficient that the reductive center of $G$ be reduced to the identity group, and that the Lie algebra of $G$ be nilpotent.
\end{remarks}

\sgasection{9}{Supplement: action of a group scheme and fixed points}
The aim of this section, added in January 2008, is to study the fixed points under the action of a group scheme.

\numpara{9.1}\textbf{Representability of the fixed-point functor.}---Let $S$ be a scheme, $G$ an $S$-group scheme acting on an $S$-scheme $X$. One defines the subfunctor $X^G$ of $X$ of fixed points of $X$ under $G$: for every $S$-scheme $T$, $X^G(T)$ is the subset of $X(T)$ formed by the fixed points under $G$ (VIII.6.e).

Let us recall the notion of an ``$S$-pure scheme'' introduced in ([G-R], §~3.3, [R]). Let $G$ be a flat $S$-group scheme of finite presentation. Since the fibers of $G$ have no embedded components, the notion of purity for $G$ takes the following form. The scheme $G$ is $S$-pure if for every strictly henselian local $S$-scheme $T$, with closed point $t$, every generic point $\widetilde x$ of a fiber of $G\times_S T$ specializes to a point of $G_t$, \ie the closure of $\widetilde x$ in $G\times_S T$ meets $G_t$.
% typo: stray quote preceding the source's opening quotation mark -> omitted.

\par\medskip\noindent\textbf{Examples.}---
\begin{enumerate}[label=(\arabic*)]
\item If $G$ is quasi-finite and separated over $S$, $G$ is $S$-pure if and only if $G$ is finite over $S$.
\item If $S=\Spec(R)$ and $G=\Spec(A)$, $G$ is $S$-pure if and only if $A$ is a projective $R$-module ([G-R], 3.3.5). In particular $G$ is $S$-pure if $G$ is diagonalizable.
\item A group scheme of multiplicative type is $S$-pure, for the notion of purity is local for the étale topology on $S$.
\item The scheme $G$ is $S$-pure if $G$ is $S$-proper or if the fibers of $G$ are irreducible or if $S$ is semilocal artinian.
\item In particular, a reductive $S$-group scheme $G$ is $S$-pure.
\end{enumerate}
\begin{proposition}{9.2}\label{XII.9.2}
Let $G$ be an $S$-group scheme of finite presentation, $S$-flat and $S$-pure, acting on an $S$-scheme $X$ of finite presentation. Then the functor $X^G$ of fixed points of $X$ under $G$ is representable by a closed subscheme of $X$, of finite presentation over $S$.
\end{proposition}
Indeed, let $a$ be in $X(S)$ and let $H$ be the group subscheme of $G$ that is the stabilizer of $a$. Then $a$ is fixed under $G$ if and only if $H=G$. Now the subfunctor $C$ of $S$ of coincidences (more precisely, if $T$ is an $S$-scheme, $C(T)=\{\varnothing\}$ if $H\times_S T\to G\times_S T$ is bijective, and $C(T)=\varnothing$ otherwise) of $H$ with $G$ is representable by a closed subscheme of $S$, defined by a sheaf of ideals of finite type ([G-R], 4.1.1). One applies this result with $S=X$, taking for $a$ the universal point of $X$ (\ie the identity in $X(X)=\Hom_S(X,X)$).

\numpara{9.3}\textbf{Infinitesimal obstruction.}---Under the hypothesis that $G$ is a flat $S$-group scheme of finite presentation acting (on the left) on a smooth $S$-prescheme $X$, we shall now study the formal smoothness of the functor $X^G$.

One supposes here that $S$ is endowed with a closed subscheme $S_0$ defined by a sheaf of quasi-coherent ideals $\mathscr I$ of square zero. One writes $j:S_0\to S$, $G_0=G\times_S S_0$, $X_0=X\times_S S_0$.
% typo: "un faisceaux" -> "un faisceau".

Let $\epsilon_0$ be a point of $X^G(S_0)$ lifting to a point of $X(S)$. We shall study the obstruction to lifting $\epsilon_0$ to a point of $X^G(S)$. Since $X$ is smooth, one knows that the liftings of $\epsilon_0$ to $X(S)$ form a trivial principal homogeneous space under the abelian group
\[
\Hom_{\OO_{S_0}}(\epsilon_0^*(\Omega^1_{X_0/S_0}),\mathscr I)
\]
(III.0.2, 0.3). One then puts
\[
L_0=\uHom_{\OO_{S_0}}(\epsilon_0^*(\Omega^1_{X_0/S_0}),\mathscr I),
\]
which is an $\OO_{S_0}$-module. Moreover, since $\epsilon_0$ is fixed under $G$, $L_0$ is naturally a $G_0$-$\OO_{S_0}$-module. One denotes by $\rho_0:G_0\to\underline{\Aut}_{S_0}(L_0)$ (resp. $\rho:G\to\underline{\Aut}_S(j_*L_0)$) the associated representation.
% typo: "pour le sous le groupe" -> "sous le groupe".
\begin{lemma}{9.4}\label{XII.9.4}
There exists a certain class $c(\epsilon_0)\in H^1(G,j_*L_0)\simeq H^1(G_0,L_0)$, defined canonically by $\epsilon_0$, such that $\epsilon_0$ lifts to $X^G(S)$ if and only if $c(\epsilon_0)=0$.
\end{lemma}
The adjunction $H^1(G,j_*L_0)\simeq H^1(G_0,L_0)$ is that of lemma III.1.1.2. We shall work with the small fppf site over $S$. For every scheme $T$ flat and of finite presentation over $S$, write $G_T=G\times_S T$, $X_T=X\times_S T$ for the corresponding objects over $T$. Consider the sheaf $A$ on the small fppf site of $S$ such that for every $T$ flat and of finite presentation over $S$, one has:
\[
A(T)=\left\{\text{set of liftings of }\epsilon_0\times_S T\text{ to }X(T)\right\}.
\]
Since the formation of $L_0$ commutes with flat base changes, the preceding considerations indicate that $A(T)$ is a trivial principal homogeneous space under the abelian group $H^0(T,j_*L_0)$. Again because $\epsilon_0$ is fixed, for every $T$ flat and of finite presentation over $S$, and every $g$ in $G(T)$, $g$ acts by affine automorphisms on $A(T)$, compatibly with the action of $G$ on $j_*L_0$. That is:
\[
g(a_T+v)=g(a_T)+\rho(g)(v),\qquad v\in H^1(T,j_*L_0).
\]
Since $G$ is flat of finite type, one can apply these considerations by taking $T=G$ and for $g$ the universal point $\mathrm{id}_G$ of $G$. One thus obtains an action of the fppf sheaf $G$ on $A$.
% typo?: H^1 rather than H^0 in the affine-action formula retained.

Let $a$ be a lifting of $\epsilon_0$ to $X(S)$. Denote by $g^\sharp$ the universal point of $G$ and define $v^\sharp\in H^0(G,j_*L_0)$ by
\[
\rho(g^\sharp)(v^\sharp)=g^\sharp\cdot a_G-a_G.
\]
For every $S$-prescheme $Y$, put $z(Y):G(Y)\to H^0(Y,j_*L_0)$, $g\mto v^\sharp(g)$. This defines the $1$-cocycle $z$ in $Z^1(G,L)$ (exposé I, §~5).
% typo?: L in Z^1(G,L) not specified, retained.

Its class $c(\epsilon_0)$ in $H^1(G,j_*L_0)$ does not depend on the choice of $a$. In particular, if $\epsilon_0$ lifts to $X^G(S)$, one has $c(\epsilon_0)=0$. Conversely, if $c(\epsilon_0)=0$, then there exists $w\in H^0(S,j_*L_0)$ such that $z(Y)(g)=g\cdot w_Y-w_Y$ for every $S$-prescheme $Y$ and every $g\in G(Y)$. Applying this to $Y=G$ and to $g^\sharp$, one concludes that $a-w\in X^G(S)$. We have therefore established that $\epsilon_0$ lifts to $X^G(S)$ if and only if $c(\epsilon_0)=0$.
% typo: missing "pour" before "tout S-préschéma" -> supplied in English.
\begin{remark}{9.5}\label{XII.9.5}
In the case where $G$ is affine and flat of finite type over affine $S$, one can refine this obstruction to a class $\widetilde c(\epsilon_0)\in H^1_{G_0}(X_0,L_0)$, where $H^1_{G_0}(X_0,L_0)$ is the $G_0$-equivariant cohomology group defined by Wevers ([W], app.~C). In this case, it is not necessary to suppose that $\epsilon_0$ lifts to $X(S)$.
% typo?: name Wevers in prose differs from Wewers in bibliography; retained.
\end{remark}

\numpara{9.6}\textbf{Smoothness of fixed points.}---
\begin{theorem}{9.7}\label{XII.9.7}
Let $G$ be a flat $S$-group scheme of finite type over a noetherian base $S$, acting on a smooth $S$-scheme $X$. Suppose that $G$ is $S$-pure, and for every geometric point $s$ over $S$ (with algebraically closed residue field), one has
\[
H^1\bigl(G_s,(\Omega^1_{X_s})^*\bigr)=0.
\]
Then the functor $X^G$ of fixed points of $X$ under $G$ is representable by a closed subscheme of $X$ smooth over $S$.
\end{theorem}
Let us prove theorem 9.7. One knows that $X^G$ is representable by a closed subscheme of $X$ of finite presentation over $S$. One can suppose $S=\Spec(A)$ local. Following the smoothness criterion (SGA~1 III.3.1.iii.bis), it suffices to verify the formal smoothness of $X^G$ for $S'=\Spec(A')$ and the closed subscheme $S'_0=\Spec(A'/\mathscr I')=\Spec(A'_0)$, where $A'$ is an artinian local ring and $\mathscr I'$ a square-zero ideal of $A'$. By reduction, one is reduced to the case where $\mathscr I'$ is annihilated by the maximal ideal of $A'$. In particular, if $k$ denotes the residue field of $A'$, $\mathscr I'$ is a $k$-vector space.

One therefore gives oneself an element $\epsilon'_0\in X^G(S'_0)$ which one wishes to lift to $X^G(S')$. One then writes $X'=X\times_S S'$, $G'=G\times_S S'$, $X'_0=X\times_S S'_0$, $G'_0=G\times_S S'_0$. Since $S'$ is affine and $X'$ is smooth over $S'$, $\epsilon'_0\in X^G(S'_0)$ lifts to $X(S')$. According to lemma 9.4, the class
\[
c(\epsilon_0)\in H^1(G'_0,L'_0)
\]
is the obstruction to lifting $\epsilon'_0$ to $X^G(S')$, where $L'_0=\uHom_{\OO_{S'_0}}(\epsilon'_0{}^*(\Omega^1_{X'_0/S'_0}),\mathscr I')$. Since $L'_0$ is a $k$-vector space, one has a canonical isomorphism $H^1(G'_0,L'_0)\isomto H^1(G'_0\times_{A'_0}k,L'_0)$. It suffices to show that $H^1(G'_0\times_{A'_0}k,L'_0)=0$. Writing $\overline k$ for an algebraic closure of $k$, one has an isomorphism (IX.3.1 in the affine case, lemma 9.11 in general)
\[
H^1(G'_0\times_{A'_0}k,L'_0)\otimes_k\overline k
\isomto H^1(G'_0\times_{A'_0}\overline k,L'_0\otimes_k\overline k).
\]
One then observes that the representation $L'_0\otimes_k\overline k$ is a direct sum of $(\Omega^1_{X'_0\times_k\overline k})^*$. By hypothesis one has $H^1(G'_0\times_{\overline k}A'_0,(\Omega^1_{X'_0\times_k\overline k})^*)=0$, which implies $H^1(G'_0\times_{A'_0}k,L'_0)=0$.
% typo?: c(epsilon_0) lacks a prime; reversed fiber-product base G'_0 times_{bar k} A'_0 and the displayed differentials retained as printed.
\begin{corollary}{9.8}\label{XII.9.8}
Let $G$ be a flat $S$-group scheme of finite presentation acting on a smooth $S$-scheme $X$ of finite presentation. Suppose that $G$ admits a composition series whose factors are of one of the following types:
\begin{enumerate}[label=(\arabic*)]
\item $S$-abelian scheme (\ie $G$ is smooth over $S$ and its fibers are abelian varieties),
\item $S$-group scheme of multiplicative type,
\item finite étale $S$-group scheme of degree invertible on $S$,
\item reductive $S$-group if $S$ is a $\QQ$-scheme.
\end{enumerate}
Then the functor $X^G$ of fixed points of $X$ under $G$ is representable by a closed subscheme of $X$, of finite presentation, smooth over $S$.
\end{corollary}
Indeed, if $1\to G'\to G\to G''\to1$ is an exact sequence of $S$-groups satisfying the hypotheses of the theorem, the $S$-group $G''$ acts on the functor $X^{G'}$ and the functor $X^G$ is nothing other than that of fixed points of $X^{G'}$ under $G''$. Thus one is reduced to the case of an $S$-group of one of the four types above. By passage to the limit, one can suppose $S$ noetherian. One will verify the cohomological criterion stated above for a fiber $G_s$ over a geometric point $s$ of $S$ and a finite-dimensional linear representation $V$ of $G_s$.

In the case of an abelian scheme, every map from $G_s^n$ to $V$ is constant and consequently the $H^i(G_s,V)$ vanish for $i>0$.

In the case of a group scheme of multiplicative type, $G_s$ is a diagonalizable group. Theorem I.5.3.3 shows that the $H^i(G_s,V)$ vanish for $i>0$.

In the next two cases, the argument is the same, since it rests on the semisimplicity of linear representations of $G_s$. Indeed, in the case of a finite étale $S$-group scheme of degree invertible on $S$, $G_s$ is a finite constant group of degree invertible in the residue field $\kappa(s)$; it is well known that every linear representation of $G_s$ is semisimple.

For the reductive case, the residue field $\kappa(s)$ is supposed (algebraically closed) of characteristic zero. One then knows that every linear representation of $G_s$ is semisimple (see [T-Y], theorem 27.3.3).
% typo: missing "que" after "On sait alors" -> supplied in English.
\begin{remark}{9.9}\label{XII.9.9}
This applies in particular to the case of an action of $G/S$ on a smooth $S$-group scheme $H$. If $G$ is a group of multiplicative type acting by inner transformations on $H$ smooth and affine over $S$, one then recovers corollary XI.5.3 asserting the smoothness of the centralizer of $G$ in $H$.
% typo: "tranformations" -> "transformations".
\end{remark}
\begin{remark}{9.10}\label{XII.9.10}
In the case where $S=\Spec(k)$ and $G$ is a linearly reductive algebraic group defined over the algebraically closed field $k$, this result is due independently to Fogarty ([F], theorem 5.4) and Iversen ([I], proposition 1.3).
\end{remark}

The proof of theorem 9.7 uses the following lemma, well known in the case of an affine group scheme (IX.3.1).
\begin{lemma}{9.11}\label{XII.9.11}
Let $G$ be a group scheme over an affine scheme $S=\Spec(A)$, and $f:S'=\Spec(A')\to S$ a flat morphism. Put $G'=G\times_S S'$. Then for every quasi-coherent $G$-$\OS$-module $M$, one has canonical isomorphisms
\[
H^i(G,M)\otimes_A A'\isomto H^i(G',M\otimes_A A')\qquad(i\ge0).
\]
\end{lemma}
Since chains of degree $n\ge1$ for Hochschild cohomology are determined by their value at the point $(\mathrm{id},\ldots,\mathrm{id})\in G(G)\times\cdots G(G)$, one has an isomorphism of $A$-modules $C^n(G,M)\simeq\Gamma(G^n,M\otimes_A\OO_{G^n})$. The cohomology group $H^n(G,M)$ is the $n$th cohomology group of the complex
\[
\begin{aligned}
L&\to\Gamma(G,M\otimes_A\OO_G)\to\Gamma(G^2,M\otimes_A\OO_{G^2})\\
&\to\Gamma(G^3,M\otimes_A\OO_{G^3})\to\cdots
\end{aligned}
\]
Since $A'$ is flat over $A$, by EGA~IV$_1$, 1.7.21, one has
\[
\Gamma(G^n,M\otimes_A\OO_{G^n})\otimes_A A'
\isomto\Gamma(G'{}^n,(M\otimes_A A')\otimes_{A'}\OO_{G'{}^n}).
\]
This implies the desired assertion on taking cohomology.
% typo?: L, rather than M, as first term of complex retained.

\begin{thebibliography}{RG71}
\item[]\nde{Additional references cited in this Exposé.}
\bibitem[Fo73]{XII.Fo73} J.~Fogarty, \textit{Fixed point schemes}, Amer.~J.~Math. \textbf{95} (1973), 35--51.
\bibitem[RG71]{XII.RG71} M.~Raynaud, L.~Gruson, \textit{Critères de platitude et de projectivité}, Invent.~math. \textbf{13} (1971), 1--89.
\bibitem[Iv72]{XII.Iv72} B.~Iversen, \textit{A fixed point formula for action of tori on algebraic varieties}, Invent.~math. \textbf{16} (1972), 229--236.
\bibitem[Ray72]{XII.Ray72} M.~Raynaud, \textit{Flat modules in algebraic geometry}, Compositio Math. \textbf{24} (1972), 11--31.
\bibitem[TY06]{XII.TY06} P.~Tauvel, R.~W.~T.~Yu, \textit{Lie algebras and algebraic groups}, Springer-Verlag, 2006.
\bibitem[We05]{XII.We05} S.~Wewers, \textit{Formal deformation of curves with group scheme action}, Ann.~Inst.~Fourier (Grenoble) \textbf{55} (2005), 1105--1165.
\end{thebibliography}

\end{document}
